% Chapter 6. All quoted numbers come from chapters/ch06-numbers.tex, written by code/ch06/run_all.py.
% Generated by code/ch06/run_all.py. Do not edit by hand.
\newcommand{\acBreachArcc}{0.02}
\newcommand{\acBreachOneHalf}{0.60}
\newcommand{\acP}{2}
\newcommand{\acQ}{22.6}
\newcommand{\acQeq}{32}
\newcommand{\acQeqTwo}{45}
\newcommand{\acQloss}{15.2}
\newcommand{\acRho}{0.2}
\newcommand{\acThrArcc}{20.6}
\newcommand{\acThrOne}{25.0}
\newcommand{\acThrOneHalf}{10.0}
\newcommand{\bdC}{79.0}
\newcommand{\bdD}{8.18}
\newcommand{\bdDate}{29 September 2026}
\newcommand{\bdDnum}{8.1797}
\newcommand{\bdDpar}{8.1797}
\newcommand{\bdDurDn}{108.18}
\newcommand{\bdDurDnTwo}{116.36}
\newcommand{\bdDurUp}{91.82}
\newcommand{\bdDurUpTwo}{83.64}
\newcommand{\bdExDn}{108.59}
\newcommand{\bdExDnTwo}{118.05}
\newcommand{\bdExUp}{92.20}
\newcommand{\bdExUpTwo}{85.12}
\newcommand{\bdFirst}{2 January 2001}
\newcommand{\bdMac}{8.34}
\newcommand{\bdMaxDurFive}{1.27}
\newcommand{\bdMaxDurTen}{2.81}
\newcommand{\bdMaxDyFive}{33}
\newcommand{\bdMaxDyTen}{26}
\newcommand{\bdMaxQdFive}{0.0081}
\newcommand{\bdMaxQdTen}{0.0255}
\newcommand{\bdNFive}{6,436}
\newcommand{\bdNTen}{6,437}
\newcommand{\bdQdDn}{108.57}
\newcommand{\bdQdDnTwo}{117.94}
\newcommand{\bdQdUp}{92.22}
\newcommand{\bdQdUpTwo}{85.22}
\newcommand{\bdRmsDurFive}{0.077}
\newcommand{\bdRmsDurTen}{0.209}
\newcommand{\bdRmsQdFive}{0.00029}
\newcommand{\bdRmsQdTen}{0.00110}
\newcommand{\bdRmsRetFive}{23.9}
\newcommand{\bdRmsRetTen}{42.0}
\newcommand{\bdY}{3.99}
\newcommand{\cdAnnCurve}{4.460}
\newcommand{\cdAnnFlat}{4.441}
\newcommand{\cdAnnHy}{3.561}
\newcommand{\cdFlatLam}{1.161}
\newcommand{\cdLamHy}{10.78}
\newcommand{\cdQHy}{650}
\newcommand{\cdUpHy}{5.34}
\newcommand{\cdUpHyOne}{19.58}
\newcommand{\cdUpIg}{1.33}
\newcommand{\cdUpIgCurve}{1.34}
\newcommand{\cdUpIgGapBp}{0.6}
\newcommand{\cpAtmFive}{3.29}
\newcommand{\cpAtmOne}{3.63}
\newcommand{\cpAtmTwo}{3.68}
\newcommand{\cpEL}{3.45}
\newcommand{\cpFFive}{4.07}
\newcommand{\cpFOne}{3.53}
\newcommand{\cpFTwo}{3.90}
\newcommand{\cpMc}{3.676}
\newcommand{\cpMcDir}{below}
\newcommand{\cpMcEq}{55.92}
\newcommand{\cpMcEqSE}{0.09}
\newcommand{\cpMcExact}{3.683}
\newcommand{\cpMcMaxEst}{23.4323}
\newcommand{\cpMcMaxName}{3 to 7\%}
\newcommand{\cpMcMaxSE}{0.0855}
\newcommand{\cpMcMaxTab}{23.6560}
\newcommand{\cpMcMaxZ}{2.6}
\newcommand{\cpMcMaxZAlt}{1.2}
\newcommand{\cpMcPaths}{400,000}
\newcommand{\cpMcSE}{0.0085}
\newcommand{\cpMcScen}{200,000}
\newcommand{\cpMcSignNeg}{5}
\newcommand{\cpMcSnB}{1.216}
\newcommand{\cpMcSnBSE}{0.019}
\newcommand{\cpMcZ}{\ensuremath{-}0.8}
\newcommand{\cpMcZabs}{0.8}
\newcommand{\cpN}{100}
\newcommand{\cpPAboveOne}{0.0}
\newcommand{\cpPAboveSix}{3.4}
\newcommand{\cpPAboveThree}{0.7}
\newcommand{\cpPAboveZero}{0.0}
\newcommand{\cpPTailOne}{3.84}
\newcommand{\cpPTailSix}{11.53}
\newcommand{\cpPTailThree}{9.74}
\newcommand{\cpPTailZero}{0.02}
\newcommand{\cpPZeroOne}{4.9}
\newcommand{\cpPZeroSix}{49.7}
\newcommand{\cpPZeroThree}{21.3}
\newcommand{\cpPZeroZero}{0.3}
\newcommand{\cpSkewHi}{5.2}
\newcommand{\cpSkewLo}{8.0}
\newcommand{\cpVolFive}{4.4}
\newcommand{\cpVolOne}{10.7}
\newcommand{\cpVolTwo}{7.2}
\newcommand{\csAboveSeventy}{29.0}
\newcommand{\csAmort}{25}
\newcommand{\csAttach}{1}
\newcommand{\csBeta}{0.5}
\newcommand{\csCap}{6}
\newcommand{\csCapZeroMax}{6.80}
\newcommand{\csCapZeroMin}{4.15}
\newcommand{\csConvMax}{12.5}
\newcommand{\csCost}{10}
\newcommand{\csCoverBetaOne}{0.992}
\newcommand{\csCoverBetaZero}{1.241}
\newcommand{\csCoverMed}{1.309}
\newcommand{\csCoverMin}{0.985}
\newcommand{\csCoverMinDate}{July 2020}
\newcommand{\csDscrBindsOne}{249}
\newcommand{\csDscrBindsZero}{249}
\newcommand{\csDscrMin}{1.25}
\newcommand{\csEl}{14.419}
\newcommand{\csEta}{77.44}
\newcommand{\csFailFirst}{July 2020}
\newcommand{\csFailLast}{September 2021}
\newcommand{\csFailYMax}{1.11}
\newcommand{\csFails}{3}
\newcommand{\csFirst}{January 2001}
\newcommand{\csG}{2}
\newcommand{\csGrade}{C}
\newcommand{\csGradeReal}{C}
\newcommand{\csLast}{September 2021}
\newcommand{\csLastB}{53.13}
\newcommand{\csLastCapOne}{5.85}
\newcommand{\csLastCapZero}{4.55}
\newcommand{\csLastCover}{0.992}
\newcommand{\csLastDscrOne}{1.24}
\newcommand{\csLastDscrZero}{1.25}
\newcommand{\csLastKOne}{7.63}
\newcommand{\csLastKZero}{5.90}
\newcommand{\csLastLoan}{61.7}
\newcommand{\csLastLtvOne}{61.7}
\newcommand{\csLastVOne}{86.04}
\newcommand{\csLastYOne}{3.69}
\newcommand{\csLastYZero}{1.11}
\newcommand{\csLoans}{100}
\newcommand{\csLtv}{65}
\newcommand{\csLtvBinds}{0}
\newcommand{\csLtvZeroMax}{61.8}
\newcommand{\csLtvZeroMin}{60.2}
\newcommand{\csMatLast}{September 2026}
\newcommand{\csMcEl}{14.501}
\newcommand{\csMcElSE}{0.112}
\newcommand{\csMcElZ}{0.7}
\newcommand{\csMcMeanFail}{55.45}
\newcommand{\csMcMeanFailSE}{0.071}
\newcommand{\csMcScen}{100,000}
\newcommand{\csMeanFail}{55.38}
\newcommand{\csMonthlyFails}{5}
\newcommand{\csN}{249}
\newcommand{\csOldBelow}{85}
\newcommand{\csOldYCut}{3.34}
\newcommand{\csPdBetaOne}{55.4}
\newcommand{\csPdBetaZero}{13.4}
\newcommand{\csPdLast}{55.4}
\newcommand{\csPdMax}{56.9}
\newcommand{\csPdMaxDate}{July 2020}
\newcommand{\csPdMed}{8.1}
\newcommand{\csPdPreMed}{7.2}
\newcommand{\csPoolTail}{89.4}
\newcommand{\csPoolTailMc}{89.4}
\newcommand{\csRateZero}{3.36}
\newcommand{\csRho}{0.308}
\newcommand{\csScore}{59}
\newcommand{\csScoreReal}{52}
\newcommand{\csScreenExitLtv}{61.8}
\newcommand{\csSenEl}{1.89}
\newcommand{\csSenElOwn}{1.91}
\newcommand{\csSenElSE}{0.013}
\newcommand{\csSenPHit}{2.7}
\newcommand{\csSenScen}{2,000,000}
\newcommand{\csSig}{18.0}
\newcommand{\csSigE}{15}
\newcommand{\csSigM}{10}
\newcommand{\csSponsor}{3}
\newcommand{\csSpread}{2.25}
\newcommand{\csYRef}{4}
\newcommand{\cvCost}{20}
\newcommand{\cvCov}{6.5}
\newcommand{\cvDrop}{23}
\newcommand{\cvDropLite}{50}
\newcommand{\cvLev}{5}
\newcommand{\cvMc}{63.8}
\newcommand{\cvMcRaw}{58.5}
\newcommand{\cvMcSe}{0.11}
\newcommand{\cvMult}{8}
\newcommand{\cvPcov}{64}
\newcommand{\cvPcovOne}{63.9}
\newcommand{\cvPlite}{21}
\newcommand{\cvRate}{10}
\newcommand{\cvRecCov}{98}
\newcommand{\cvRecLite}{64}
\newcommand{\cvSig}{25}
\newcommand{\cvT}{5}
\newcommand{\cvTimes}{3}
\newcommand{\dlBaseHi}{1.0}
\newcommand{\dlBaseLo}{1.5}
\newcommand{\dlBaseNaArcc}{3.12}
\newcommand{\dlBaseNaMax}{4.00}
\newcommand{\dlBetaOne}{0.06}
\newcommand{\dlBetaOneT}{2.2}
\newcommand{\dlBetaSum}{0.10}
\newcommand{\dlBetaZero}{0.05}
\newcommand{\dlBillFirst}{3.70}
\newcommand{\dlBillHighQ}{2023Q4}
\newcommand{\dlBillLast}{3.7}
\newcommand{\dlBillPeak}{5.5}
\newcommand{\dlCap}{5}
\newcommand{\dlCatchQ}{2.1875}
\newcommand{\dlCoverArcc}{189}
\newcommand{\dlCoverMin}{150}
\newcommand{\dlDeArcc}{1.12}
\newcommand{\dlDeMax}{2}
\newcommand{\dlDimN}{20}
\newcommand{\dlEarly}{98}
\newcommand{\dlExBase}{3.12}
\newcommand{\dlExFeeShare}{37}
\newcommand{\dlExGross}{15.1}
\newcommand{\dlExInc}{2.40}
\newcommand{\dlExNet}{9.6}
\newcommand{\dlExRd}{6}
\newcommand{\dlFeeNa}{5.3}
\newcommand{\dlFeeShare}{22}
\newcommand{\dlFillTen}{50}
\newcommand{\dlFillTwenty}{25}
\newcommand{\dlFirms}{19}
\newcommand{\dlFirst}{2021Q2}
\newcommand{\dlHurdleA}{7}
\newcommand{\dlHurdleQ}{1.75}
\newcommand{\dlIcrLev}{2.7}
\newcommand{\dlIcrPrivate}{2.0}
\newcommand{\dlIncRate}{20}
\newcommand{\dlLast}{2026Q2}
\newcommand{\dlLeftFourTen}{6}
\newcommand{\dlLeftFourTwenty}{32}
\newcommand{\dlLevMax}{1.67}
\newcommand{\dlLevMed}{1.26}
\newcommand{\dlLgdDirect}{67}
\newcommand{\dlLgdHy}{61}
\newcommand{\dlLgdSynd}{48}
\newcommand{\dlLossDirectTwo}{1.34}
\newcommand{\dlLossSyndTwo}{0.96}
\newcommand{\dlMarkFirst}{103.2}
\newcommand{\dlMarkLast}{97.9}
\newcommand{\dlMarkMin}{97.9}
\newcommand{\dlMarkMinQ}{2026Q2}
\newcommand{\dlNonaccCost}{1.8}
\newcommand{\dlNonaccFv}{1.2}
\newcommand{\dlPass}{1.14}
\newcommand{\dlPassRsq}{0.91}
\newcommand{\dlPassSe}{0.10}
\newcommand{\dlPdPerDirect}{1.49}
\newcommand{\dlPdPerHy}{1.64}
\newcommand{\dlPdPerSynd}{2.08}
\newcommand{\dlPfMin}{\ensuremath{-}7.3}
\newcommand{\dlPfMinDs}{\ensuremath{-}8.7}
\newcommand{\dlPfMinQ}{2020Q1}
\newcommand{\dlPfN}{54}
\newcommand{\dlPfPhi}{0.12}
\newcommand{\dlPfRatio}{1.13}
\newcommand{\dlPfVol}{4.5}
\newcommand{\dlPfVolDs}{5.1}
\newcommand{\dlQuartersTen}{3.3}
\newcommand{\dlQuartersTwenty}{8.0}
\newcommand{\dlRecDirect}{33}
\newcommand{\dlRecHy}{39}
\newcommand{\dlRecSynd}{52}
\newcommand{\dlRho}{0.07}
\newcommand{\dlRhoEw}{0.36}
\newcommand{\dlRhoN}{350}
\newcommand{\dlRhoSe}{0.05}
\newcommand{\dlRoeMean}{8.8}
\newcommand{\dlRoeMin}{\ensuremath{-}0.04}
\newcommand{\dlRoeMinQ}{2022Q2}
\newcommand{\dlRoeN}{20}
\newcommand{\dlRoeVol}{2.0}
\newcommand{\dlSpreadFirstBp}{727}
\newcommand{\dlSpreadLastBp}{682}
\newcommand{\dlYieldArcc}{10.3}
\newcommand{\dlYieldFirst}{11.0}
\newcommand{\dlYieldFirstQ}{2022Q4}
\newcommand{\dlYieldLast}{10.5}
\newcommand{\dlYieldN}{15}
\newcommand{\dlYieldPeak}{12.7}
\newcommand{\dlYieldPeakQ}{2023Q4}
\newcommand{\hlAdjFive}{6.0}
\newcommand{\hlAdjOne}{0.3}
\newcommand{\hlAdjTen}{23.3}
\newcommand{\hlAdjTwo}{1.0}
\newcommand{\hzAvgFive}{1.185}
\newcommand{\hzFlatPar}{60.28}
\newcommand{\hzFlatTri}{60.00}
\newcommand{\hzLFive}{1.585}
\newcommand{\hzLOne}{0.664}
\newcommand{\hzLSeven}{1.830}
\newcommand{\hzLTen}{2.010}
\newcommand{\hzLThree}{1.045}
\newcommand{\hzMcN}{1,000,000}
\newcommand{\hzMcPd}{5.765}
\newcommand{\hzMcPdSE}{0.023}
\newcommand{\hzMcPdZ}{0.5}
\newcommand{\hzMcProt}{312.84}
\newcommand{\hzMcProtSE}{1.27}
\newcommand{\hzMcProtZ}{0.5}
\newcommand{\hzPdFive}{5.75}
\newcommand{\hzPdFiveX}{5.753}
\newcommand{\hzProt}{312.19}
\newcommand{\hzQFive}{70}
\newcommand{\hzQOne}{40}
\newcommand{\hzQSeven}{80}
\newcommand{\hzQTen}{90}
\newcommand{\hzQThree}{55}
\newcommand{\hzRec}{40}
\newcommand{\hzSFive}{94.25}
\newcommand{\hzTriFive}{1.167}
\newcommand{\hzTriOne}{0.667}
\newcommand{\hzTriSeven}{1.333}
\newcommand{\hzTriTen}{1.500}
\newcommand{\hzTriThree}{0.917}
\newcommand{\icBillHigh}{5.5}
\newcommand{\icBillLow}{0.04}
\newcommand{\icCostRise}{80}
\newcommand{\icLow}{3.6}
\newcommand{\imAone}{51.66}
\newcommand{\imAtwo}{50.89}
\newcommand{\imBarDn}{0.170}
\newcommand{\imBarDnTwo}{0.736}
\newcommand{\imBarUp}{0.145}
\newcommand{\imBarUpTwo}{0.534}
\newcommand{\imShoDn}{\ensuremath{-}4.09}
\newcommand{\imShoDnTwo}{\ensuremath{-}8.57}
\newcommand{\imShoLoss}{4.09}
\newcommand{\imShoUp}{3.74}
\newcommand{\imShoUpTwo}{7.16}
\newcommand{\imSingle}{83.95}
\newcommand{\imZ}{3.5}
\newcommand{\lhEqOne}{79.18}
\newcommand{\lhEqSix}{36.66}
\newcommand{\lhEqThree}{58.52}
\newcommand{\lhMzAOne}{22.81}
\newcommand{\lhMzASix}{18.65}
\newcommand{\lhMzAThree}{23.18}
\newcommand{\lhMzBOne}{4.22}
\newcommand{\lhMzBSix}{12.40}
\newcommand{\lhMzBThree}{10.96}
\newcommand{\lhSnAOne}{0.69}
\newcommand{\lhSnASix}{8.60}
\newcommand{\lhSnAThree}{5.28}
\newcommand{\lhSnBOne}{0.02}
\newcommand{\lhSnBSix}{4.03}
\newcommand{\lhSnBThree}{1.12}
\newcommand{\lhSupOne}{0.0000}
\newcommand{\lhSupSix}{0.29}
\newcommand{\lhSupThree}{0.01}
\newcommand{\mcDD}{0.50}
\newcommand{\mcLev}{0.697}
\newcommand{\mcMu}{7}
\newcommand{\mcPdP}{30.8}
\newcommand{\mcPdQ}{41.9}
\newcommand{\mcSigV}{30.0}
\newcommand{\mcSpread}{314}
\newcommand{\mcV}{98.85}
\newcommand{\mtPeakNine}{642}
\newcommand{\mtPeakSeven}{221}
\newcommand{\mtPeakTNine}{0.5}
\newcommand{\mtPeakTSeven}{6.4}
\newcommand{\mtR}{3}
\newcommand{\mtSNineFive}{398.9}
\newcommand{\mtSNineOne}{603.6}
\newcommand{\mtSNineQ}{293.2}
\newcommand{\mtSNineTen}{316.9}
\newcommand{\mtSSevenFive}{218.8}
\newcommand{\mtSSevenOne}{103.2}
\newcommand{\mtSSevenQ}{\ensuremath{7.5\times 10^{-4}}}
\newcommand{\mtSSevenTen}{215.6}
\newcommand{\mtSigV}{25}
\newcommand{\oiApproxFive}{40}
\newcommand{\oiApproxThree}{67}
\newcommand{\oiBpFive}{50}
\newcommand{\oiBpThree}{77}
\newcommand{\oiC}{9}
\newcommand{\oiFirstFive}{50}
\newcommand{\oiFirstThree}{76}
\newcommand{\oiU}{2}
\newcommand{\pcCalCap}{7.00}
\newcommand{\pcCalConst}{9.50}
\newcommand{\pcCalDscr}{1.23}
\newcommand{\pcCalDy}{11.67}
\newcommand{\pcCalExitLtv}{59.7}
\newcommand{\pcCalGrade}{B}
\newcommand{\pcCalLtv}{60.0}
\newcommand{\pcCalScore}{64}
\newcommand{\pcOttCap}{7.17}
\newcommand{\pcOttConst}{11.12}
\newcommand{\pcOttDscr}{0.92}
\newcommand{\pcOttDy}{10.24}
\newcommand{\pcOttExitLtv}{66.6}
\newcommand{\pcOttGrade}{D}
\newcommand{\pcOttLtv}{70.0}
\newcommand{\pcOttScore}{38}
\newcommand{\pcTorCap}{4.86}
\newcommand{\pcTorConst}{11.50}
\newcommand{\pcTorDscr}{0.56}
\newcommand{\pcTorDy}{6.43}
\newcommand{\pcTorExitLtv}{89.4}
\newcommand{\pcTorGrade}{D}
\newcommand{\pcTorLtv}{75.7}
\newcommand{\pcTorScore}{13}
\newcommand{\slAmort}{25}
\newcommand{\slBm}{67.45}
\newcommand{\slBs}{53.96}
\newcommand{\slCost}{10}
\newcommand{\slG}{2}
\newcommand{\slMezFifteen}{9.95}
\newcommand{\slMezRatio}{8.8}
\newcommand{\slMezTen}{2.16}
\newcommand{\slMezTwenty}{19.11}
\newcommand{\slPdMezFifteen}{16.2}
\newcommand{\slPdMezTen}{5.3}
\newcommand{\slPdMezTwenty}{26.0}
\newcommand{\slPdSenFifteen}{4.92}
\newcommand{\slPdSenTen}{0.44}
\newcommand{\slPdSenTwenty}{12.68}
\newcommand{\slRate}{6}
\newcommand{\slSenFifteen}{0.610}
\newcommand{\slSenRatio}{82}
\newcommand{\slSenTen}{0.029}
\newcommand{\slSenTwenty}{2.359}
\newcommand{\slT}{5}
\newcommand{\srCirInvKap}{0.258}
\newcommand{\srCirInvKapx}{0.25792}
\newcommand{\srCirInvMax}{17.1}
\newcommand{\srCirInvR}{5.22}
\newcommand{\srCirInvRbar}{2.29}
\newcommand{\srCirInvRbarx}{2.28561}
\newcommand{\srCirInvRms}{9.8}
\newcommand{\srCirInvRx}{5.21764}
\newcommand{\srCirNowKap}{1.103}
\newcommand{\srCirNowKapx}{1.10274}
\newcommand{\srCirNowMax}{9.2}
\newcommand{\srCirNowR}{2.16}
\newcommand{\srCirNowRbar}{4.06}
\newcommand{\srCirNowRbarx}{4.06386}
\newcommand{\srCirNowRms}{5.0}
\newcommand{\srCirNowRx}{2.16297}
\newcommand{\srDateInv}{31 July 2023}
\newcommand{\srDateNow}{29 September 2026}
\newcommand{\srFeller}{38.3}
\newcommand{\srGapInv}{0.088}
\newcommand{\srGapNow}{0.003}
\newcommand{\srHalfNow}{7.5}
\newcommand{\srMcExact}{0.677851}
\newcommand{\srMcMonth}{0.677941}
\newcommand{\srMcMonthDir}{above}
\newcommand{\srMcMonthSE}{0.000054}
\newcommand{\srMcMonthZ}{1.7}
\newcommand{\srMcMonthZabs}{1.7}
\newcommand{\srMcPaths}{100,000}
\newcommand{\srMcWeek}{0.677808}
\newcommand{\srMcWeekDir}{below}
\newcommand{\srMcWeekSE}{0.000054}
\newcommand{\srMcWeekZ}{\ensuremath{-}0.8}
\newcommand{\srMcWeekZabs}{0.8}
\newcommand{\srMeanBill}{1.947}
\newcommand{\srNegOne}{0.00}
\newcommand{\srNegOneProb}{\ensuremath{1.2\times 10^{-14}}}
\newcommand{\srNegTen}{0.40}
\newcommand{\srSigC}{0.04837}
\newcommand{\srSigV}{0.6749}
\newcommand{\srVasInvKap}{0.258}
\newcommand{\srVasInvKapx}{0.25796}
\newcommand{\srVasInvMax}{17.2}
\newcommand{\srVasInvR}{5.22}
\newcommand{\srVasInvRbar}{2.26}
\newcommand{\srVasInvRbarx}{2.26233}
\newcommand{\srVasInvRms}{9.9}
\newcommand{\srVasInvRx}{5.21914}
\newcommand{\srVasNowKap}{1.104}
\newcommand{\srVasNowKapx}{1.10421}
\newcommand{\srVasNowMax}{9.2}
\newcommand{\srVasNowR}{2.16}
\newcommand{\srVasNowRbar}{4.06}
\newcommand{\srVasNowRbarx}{4.06190}
\newcommand{\srVasNowRms}{5.0}
\newcommand{\srVasNowRx}{2.16291}
\newcommand{\srVolN}{309}
\newcommand{\swAnnFiveFive}{3.7197}
\newcommand{\swFive}{3.732}
\newcommand{\swFiveFive}{4.100}
\newcommand{\swTen}{3.898}
\newcommand{\trEqOne}{74.30}
\newcommand{\trEqSix}{35.52}
\newcommand{\trEqThree}{56.04}
\newcommand{\trMzAOne}{24.63}
\newcommand{\trMzASix}{18.76}
\newcommand{\trMzAThree}{23.66}
\newcommand{\trMzBOne}{5.78}
\newcommand{\trMzBSix}{12.55}
\newcommand{\trMzBThree}{11.51}
\newcommand{\trSnAOne}{1.15}
\newcommand{\trSnASix}{8.72}
\newcommand{\trSnAThree}{5.65}
\newcommand{\trSnBOne}{0.04}
\newcommand{\trSnBSix}{4.10}
\newcommand{\trSnBThree}{1.24}
\newcommand{\trSupOne}{0.0000}
\newcommand{\trSupSix}{0.30}
\newcommand{\trSupThree}{0.02}
\newcommand{\utCost}{20}
\newcommand{\utElF}{2.3}
\newcommand{\utElL}{25.2}
\newcommand{\utElU}{16.0}
\newcommand{\utFace}{50}
\newcommand{\utFo}{20}
\newcommand{\utLo}{30}
\newcommand{\utLsBlend}{366}
\newcommand{\utLsF}{46}
\newcommand{\utLsL}{580}
\newcommand{\utLsU}{349}
\newcommand{\utPi}{316}
\newcommand{\utSf}{362}
\newcommand{\utSig}{35}
\newcommand{\utSl}{896}
\newcommand{\utSu}{682}
\newcommand{\utT}{5}
\newcommand{\utV}{100}
\newcommand{\utW}{40}
\newcommand{\xdA}{\ensuremath{-}0.919}
\newcommand{\xdBreach}{35.8}
\newcommand{\xdEq}{22.5}
\newcommand{\xdFill}{40}
\newcommand{\xdIcr}{1.91}
\newcommand{\xdLeft}{1.7}
\newcommand{\xdPd}{2.99}
\newcommand{\xdSl}{767}
\newcommand{\xdThr}{16.7}
\newcommand{\xsEightDscr}{1.224}
\newcommand{\xsEightDy}{8.571}
\newcommand{\xsEightLtv}{70}
\newcommand{\xsFiveDebt}{86.964}
\newcommand{\xsFiveE}{33.036}
\newcommand{\xsFivePd}{27.00}
\newcommand{\xsFiveS}{165.6}
\newcommand{\xsFour}{4.875}
\newcommand{\xsNineB}{871,046}
\newcommand{\xsNineCover}{1.069}
\newcommand{\xsNineK}{7.732}
\newcommand{\xsNineMax}{931,241}
\newcommand{\xsOneC}{23.194}
\newcommand{\xsOneD}{4.4581}
\newcommand{\xsOneDur}{93.4925}
\newcommand{\xsOneEx}{93.5199}
\newcommand{\xsOneP}{95.6240}
\newcommand{\xsOneQd}{93.5202}
\newcommand{\xsSeven}{12.86}
\newcommand{\xsSixL}{2.50}
\newcommand{\xsSixS}{88.25}
\newcommand{\xsThreeA}{18.194}
\newcommand{\xsThreeB}{32.598}
\newcommand{\xsThreePV}{39.331}
\newcommand{\xsTwoD}{1.903864}
\newcommand{\xsTwoNum}{1.903864}

\begin{chapterrefs}
\setcounter{chapter}{5}
\chapter{Interest Rates, Credit Risk, Private Credit and Structured Products}
\label{ch:credit}

A private credit fund lends to a company or against a building, collects a coupon above the
government rate, and hopes to be repaid. Whether it is repaid depends largely on two things. The first is the level of
interest rates when the loan has to be refinanced, which no one knows at origination. The second is
whether the borrower's income and assets hold up, which depends partly on the borrower and partly on
the economy that every other borrower in the fund shares. This chapter builds the mathematics for
both and then puts them together on a question with Canadian data: could a five year commercial
mortgage originated in any month from \csFirst{} through \csLast{} have been refinanced when it came due?

The rate half starts with a single bond. Duration and convexity say how its price responds to a
change in yield, and the same algebra tells a lender how to fund a loan so that a parallel move in
rates cannot hurt. Models of the short rate then combine today's curve with assumptions about how rates move to give a
whole distribution of future curves, which is what a cap or a refinancing test needs; a swap, it turns out, needs only today's curve. The credit
half starts with a firm whose debt is a short put on its assets, moves to a description of default
as a sudden event with an intensity that can be read from credit default swaps, and ends with pools:
how a common factor makes defaults arrive together, and how a structure cuts the pool's losses into
tranches. The private credit section asks what a lender's screening ratios measure and how they
relate to the expected loss of a slice of the capital stack. The case study joins the two halves.

\begin{keybox}[How to read this chapter]
The numbers in the text are computed by nine scripts in \texttt{code/ch06/}
(\texttt{bonds\_duration}, \texttt{short\_rates}, \texttt{swaps\_caps}, \texttt{merton\_model},
\texttt{hazard\_cds}, \texttt{copula\_tranches}, \texttt{private\_credit}, \texttt{case\_refinance} and
\texttt{exercises}); the tests in \texttt{code/ch06/test\_*.py} check them. An independent R script
recomputes the duration tests, refits the curve, prices caplets, calibrates the Merton model,
bootstraps the hazard curve, evaluates the copula by its own quadrature, rescores the screener's
example loans and rebuilds the case study from the data file. GNU Octave solves the CIR equations
numerically and checks the caplet, the calibration and the large-pool losses. The screener's own
JavaScript is run by the tests to confirm that the Python port gives the same grades.

Random draws come from NumPy's default generator with base seed \texttt{606} and a stream number
fixed in each script. The seed is an arbitrary label that makes the runs repeatable; it carries no
date or other meaning. Each Monte Carlo standard error is the sample standard deviation of the
quantity actually averaged, divided by the square root of the number of draws.

Rates follow one convention per instrument, and the conversion is stated wherever two meet. Bond
yields $y$ and coupon rates $c$ are compounded twice a year, as Government of Canada bonds pay.
Treasury bill yields are simple rates on an actual/365 basis. Zero-coupon prices are $P(t,T)$, and
rates inside the short-rate models are continuously compounded. Forward and swap rates on a
quarterly or semi-annual schedule are simple. Loan rates in the private credit sections are annual
rates compounded monthly, the convention of the screener used there. Canadian mortgages repaid in
blended payments must state their rate calculated yearly or half-yearly under section 6 of the
Interest Act, so a quoted 6\% compounded half-yearly is 6.09\% effective a year, while 6\% compounded
monthly is 6.17\% effective. The stylised loans of \cref{sec6:private} keep the screener's monthly
convention; the case study of \cref{sec6:case} compounds half-yearly, as a Canadian mortgage would.

A few letters change meaning between sections. $D$ and $C$ are a bond's modified duration and
convexity. $r_t$ is the short rate in \cref{sec6:short,sec6:swaps}; in \cref{sec6:merton} $r$ is a
constant rate. $A(x)$ and $B(x)$ are the coefficients of a bond price with time to maturity $x$.
$K$ is a strike, which in the Merton model is the face value of the debt. $V$ is the value of the
assets behind a claim, a firm's in \cref{sec6:merton} and a property's in
\cref{sec6:private,sec6:case}. $\tau$ is always a default time. $\rho$ is the correlation of the
common factor in a pool. Volatilities are $\sigma$ for the short rate (and, in \cref{sec6:case}, for the total income shock)
and $\sigma$ with a subscript
otherwise ($\sigma_V$ for assets, $\sigma_M$ and $\sigma_\varepsilon$ for income shocks). $\mathbb P$
is the real-world probability, $Q$ the risk-neutral one and $Q^T$ the forward measure for maturity $T$.
$L$ is a liability in \cref{sec6:duration}, a simple rate when it has two arguments, and a pool's loss
in \cref{sec6:copula}; $R$ with two arguments is a rate and alone a recovery rate. $S(t)$ is a
survival probability and $S_{\mathrm{sw}}$ a swap rate. $p(y)$ is a bond price in
\cref{sec6:duration}, and elsewhere $p$ is a default probability. $X_j$ are asset payments in
\cref{sec6:duration} and $X_i$ the latent variables of a pool. $s$ is a credit spread, $\lambda$ a
hazard rate, $\ell$ a quasi-leverage, $\gamma$ a refinancing cover, $\zeta$ an income shock and
$\sigma_{\mathrm B}$ a Black volatility; $d_1$ and $d_2$ are defined where they are used. $\mathcal L$ is a loan amount and $\mathcal N$ a
property's net operating income, not to be confused with $L$ above or the default count $N$.
\end{keybox}

% =====================================================================================
\section{How far does a bond's price move when its yield moves?}
\label{sec6:duration}

A bond with coupon\ix{Coupon} rate $c$, face value 100 and maturity $T$ years on a coupon date pays $50c$ at
each half year $t_i=i/2$, $i=1,\dots,2T$, and 100 more at $T$. Write $\mathrm{CF}_i$ for the payment at
$t_i$. At yield $y$, compounded twice a year, its price is
\begin{equation}\label{eq6:price}
  p(y)=\sum_{i=1}^{2T}\mathrm{CF}_i\Bigl(1+\frac y2\Bigr)^{-i}.
\end{equation}
The price falls as the yield rises, and the question for anyone holding the bond is by how much.

\begin{definition}\label{def6:duration}
The \emph{modified duration\ix{Duration}\ix{Duration!modified}}\index{Duration!modified} and \emph{convexity\ix{Convexity}} of the bond at yield $y$ are
\[
  D=-\frac{p'(y)}{p(y)},\qquad C=\frac{p''(y)}{p(y)} .
\]
The \emph{Macaulay\ix{Duration!Macaulay} duration}\index{Macaulay duration} is $D_{\mathrm{mac}}=\sum_it_i\,\mathrm{CF}_i(1+y/2)^{-i}/p(y)$, the
present-value weighted average time of the payments, and $D=D_{\mathrm{mac}}/(1+y/2)$.
\end{definition}

The last identity follows by differentiating \cref{eq6:price} term by term, since
$\tfrac{d}{dy}(1+y/2)^{-i}=-t_i(1+y/2)^{-i-1}$ with $t_i=i/2$. A second-order Taylor expansion gives the
familiar approximation for a yield change $\Delta y$,
\begin{equation}\label{eq6:taylor}
  \frac{p(y+\Delta y)-p(y)}{p(y)}\approx-D\,\Delta y+\tfrac12C\,(\Delta y)^2 .
\end{equation}
Chapter~1 used exactly this to turn daily yield changes into returns. The sign of the error is
known.

\begin{proposition}[Price and yield]\label{prop6:shape}
For a bond with nonnegative payments, not all zero, and any $y>-2$:
\begin{enumerate}[label=(\roman*)]
  \item $p'(y)<0<p''(y)$ and $p'''(y)<0$;
  \item the error of \cref{eq6:taylor}, actual minus approximate relative price change\ix{Price change}, has the sign of
  $-\Delta y$: the approximation overstates the price after a rise in yield and understates it after a
  fall.
\end{enumerate}
\end{proposition}

\begin{proof}
(i) Write $v=1+y/2>0$. The $k$th derivative of $v^{-i}$ in $y$ is
$(-1)^k\,i(i+1)\cdots(i+k-1)\,2^{-k}v^{-i-k}$, so every term of $p^{(k)}$ has the sign of $(-1)^k$ and at
least one term is nonzero. (ii) By Taylor's theorem with the Lagrange remainder, for some $\xi$ between
$y$ and $y+\Delta y$, $p(y+\Delta y)=p(y)+p'(y)\Delta y+\tfrac12p''(y)\Delta y^2+\tfrac16p'''(\xi)\Delta y^3$.
Dividing by $p(y)>0$, the error is $p'''(\xi)\Delta y^3/(6p(y))$, whose sign is that of $-\Delta y$ by (i).
\end{proof}

For a bond priced at par the duration has a closed form, which is handy for checking code.

\begin{proposition}[Duration of a par bond]\label{prop6:par}
If $c=y>0$, so that $p(y)=100$, then $D=\bigl(1-(1+y/2)^{-2T}\bigr)/y$.
\end{proposition}

\begin{proof}
For fixed $c$, summing the geometric series in \cref{eq6:price} with $v=1+y/2$ gives
$p(y)=100\,c/y\,(1-v^{-2T})+100\,v^{-2T}$. Differentiate in $y$ with $c$ held fixed, using
$dv^{-2T}/dy=-Tv^{-2T-1}$:
\[
  p'(y)=-\frac{100c}{y^2}(1-v^{-2T})+\frac{100c}{y}\,T\,v^{-2T-1}-100\,T\,v^{-2T-1}.
\]
At $c=y$ the last two terms cancel, leaving $p'(y)=-100(1-v^{-2T})/y$, and $p(y)=100$.
\end{proof}

\begin{example}{A ten year Government of Canada bond}\label{ex6:tenyear}
On \bdDate{} the ten year benchmark\ix{Benchmark} yield was \bdY\%. Treat it as the yield of a par bond\ix{Par bond} with
semi-annual coupons, the same idealisation used for the curve in Chapter~4. Its modified duration
is \bdD{} years (\cref{prop6:par} and the sum in \cref{def6:duration} both give \bdDpar), its Macaulay duration \bdMac{} years and its convexity \bdC. \Cref{tab6:approx}
compares the exact price after a parallel move of one or two percentage points\ix{Percentage point} with the two
approximations. Duration alone misses by several tenths of a point of price for a one point move and
by more than a point for two; with convexity the misses shrink to hundredths and about a tenth, with
the signs that \cref{prop6:shape} predicts.
\end{example}

\begin{table}[t]
\centering
\caption{\kind{exact} Price per 100 face of the ten year par bond of \cref{ex6:tenyear} after a parallel move in
its yield from \bdY\%.}
\label{tab6:approx}
\begin{tabular}{lrrrr}
\toprule
yield change & $-2$ points & $-1$ point & $+1$ point & $+2$ points\\
\midrule
exact, \cref{eq6:price}         & \bdExDnTwo & \bdExDn & \bdExUp & \bdExUpTwo\\
duration only                   & \bdDurDnTwo & \bdDurDn & \bdDurUp & \bdDurUpTwo\\
duration and convexity          & \bdQdDnTwo & \bdQdDn & \bdQdUp & \bdQdUpTwo\\
\bottomrule
\end{tabular}
\end{table}

\begin{figure}[t]
  \centering
  \includegraphics[width=0.84\textwidth]{ch06/fig01_price_yield}
  \caption{\kind{exact} Price of the ten year par bond of \cref{ex6:tenyear} against its yield, with the tangent
  line of the duration approximation and the parabola of \cref{eq6:taylor}, both taken at \bdY\%.
  The parabola lies above the price curve to the right of the dot and below it to the left.}
  \label{fig6:priceyield}
\end{figure}

\paragraph{Is convexity worth the trouble for daily moves?}
For daily changes of the size seen in Canadian benchmark yields\ix{Benchmark bond}, convexity adds little to the price
approximation. The file of benchmark yields
used in Chapter~4 runs from \bdFirst{} to \bdDate. For each pair of consecutive observations, take a
par bond at the first yield, reprice it exactly at the second, and compare with the two
approximations. This tests the approximation along observed yield changes; it is a statement about
the pricing function and makes no claim about what any bond earned. Over \bdNTen{} pairs for the ten
year bond, the exact price changes have a root mean square\ix{Mean squared error} of \bdRmsRetTen{} basis points\ix{Basis point} of price. The
duration-only error has a root mean square of \bdRmsDurTen{} and a maximum of \bdMaxDurTen{} basis points
of price; the largest daily move in the ten year yield was \bdMaxDyTen{} basis points. Adding convexity brings the
root mean square down to \bdRmsQdTen{} and the maximum to \bdMaxQdTen. For the five year bond the
figures are \bdRmsDurFive{} and \bdMaxDurFive{} with duration alone and \bdRmsQdFive{} and \bdMaxQdFive{}
with convexity. On every day with a nonzero change, the sign of the remaining error was that of
\cref{prop6:shape}(ii). \Cref{fig6:daily} shows the pattern: the duration error is the parabola
$\tfrac12C\Delta y^2$, and what is left after convexity is the cubic term.

\begin{figure}[t]
  \centering
  \includegraphics[width=0.94\textwidth]{ch06/fig02_daily_error}
  \caption{\kind{historical} Error of the two approximations in \cref{eq6:taylor} for one-day changes in the five and ten year
  benchmark yields, \bdFirst{} to \bdDate, in basis points of price. Each point reprices a par bond at
  the next day's yield. The error after convexity is drawn 100 times larger so that it shows.}
  \label{fig6:daily}
\end{figure}

\subsection{Funding a liability so that a parallel move cannot hurt}

A lender who has promised a payment, or an insurer who owes one, wants assets whose value moves
with the obligation. Suppose for this subsection that the curve is flat at a continuously compounded
rate $z$ and moves only in parallel, to $z+\Delta z$. For cash flows $X_j\ge0$ at times $s_j$ the
present value\ix{Present value} is $\sum_jX_je^{-zs_j}$ and the duration\ix{Duration}, the present-value weighted average time, is
also minus the derivative of the log value in $z$. A portfolio's duration is therefore the
value-weighted average of its parts' durations, because the derivative of a sum is the sum of the
derivatives.

\begin{theorem}[Immunisation against parallel shifts]\label{thm6:redington}
Let a liability $L>0$ be due at $T$, and let assets pay $X_j\ge0$ at times $s_j$. Suppose that at the
rate $z$ the assets and the liability have the same present value and the same duration,
\[
  \sum_jX_je^{-zs_j}=Le^{-zT},\qquad \sum_js_jX_je^{-zs_j}=T\,Le^{-zT}.
\]
Then for every shift $\Delta z\in\mathbb R$ the surplus
$\Sigma(\Delta z)=\sum_jX_je^{-(z+\Delta z)s_j}-Le^{-(z+\Delta z)T}$ is nonnegative, and it is zero only at
$\Delta z=0$ unless every payment with $X_j>0$ falls at $T$.
\end{theorem}

\begin{proof}
Put $w_j=X_je^{-zs_j}\ge0$, so that $\sum_jw_j=Le^{-zT}$ and $\sum_jw_js_j=TLe^{-zT}$. Then
\[
  e^{(z+\Delta z)T}\,\Sigma(\Delta z)=\sum_jw_je^{-\Delta z(s_j-T)}-Le^{-zT}
  \ \ge\ \sum_jw_j\bigl(1-\Delta z(s_j-T)\bigr)-Le^{-zT}=0,
\]
using $e^u\ge1+u$, with equality only at $u=0$. Equality for some $\Delta z\ne0$ would need
$\Delta z(s_j-T)=0$ for every $j$ with $w_j>0$, that is $s_j=T$.
\end{proof}

The result goes back to \textcite{redington1952}, who stated the local version with convexity, and
\textcite{fisher1971}, who studied bond portfolios held to a horizon. Its hypotheses do the work. With
one flat rate and parallel moves there is a single risk factor; when the short and long ends move
differently a duration-matched book can lose, and the theorem says nothing about that. The asset
payments must be nonnegative, so a position funded partly with short bonds is not covered.
\Cref{fig6:immunise} shows the two cases for 100 owed in seven years at $z=\imZ\%$. Zero-coupon
bonds\ix{Zero-coupon bond} paying \imAone{} at two years and \imAtwo{} at fifteen years match both present value and
duration, and the surplus stays positive: \imBarUp{} after a one point rise and \imBarDn{} after a one
point fall. A single two year zero paying \imSingle{} matches only the present value; it gains
\imShoUp{} if rates rise by one point and loses \imShoLoss{} if they fall. A two year asset and a seven year liability respond differently to a
parallel shift of the discount rate: when rates fall the liability gains more value than the asset,
whose proceeds must then be reinvested for five years at the lower rate. That is the funding problem of a lender whose loans are
longer than its liabilities, seen from the other side, and it returns in \cref{sec6:structures} and
in the case study.

\begin{figure}[t]
  \centering
  \includegraphics[width=0.84\textwidth]{ch06/fig03_immunise}
  \caption{\kind{exact} Surplus of assets over a liability of 100 due in seven years after a parallel shift of a
  flat continuously compounded curve at \imZ\%. The barbell of two and fifteen year zeros matches the
  liability's present value and duration, as in \cref{thm6:redington}; the single two year zero
  matches the present value only.}
  \label{fig6:immunise}
\end{figure}

% =====================================================================================
\section{What does a model of the short rate add?}
\label{sec6:short}

A curve on one date prices any fixed set of cash flows by discounting. It cannot by itself price a payment whose amount depends on future rates, such as a cap or an option to
prepay\ix{Prepayment}, nor say how likely a loan is to be refinanceable in five years. Pricing needs a model of how the
curve moves under the pricing measure; the refinancing question needs one under the real-world measure,
and the two need not agree, which is why the case study in \cref{sec6:case} uses the rates that actually
occurred. The simplest models
describe a single number, the instantaneous short rate\ix{Short rate} $r_t$, under a risk-neutral measure\ix{Equivalent martingale measure} $Q$ in
which the price at $t$ of a payment of 1 at $T$ is
\begin{equation}\label{eq6:bond}
  P(t,T)=\E_Q\Bigl[e^{-\int_t^Tr_u\,du}\Bigm|\mathcal F_t\Bigr].
\end{equation}
The two models below were met in Chapters 3 and 4 as processes; here they produce bond prices\ix{Bond price}.

\begin{proposition}[Vasicek\index{Vasicek model} bond prices]\label{prop6:vasicek}
Under $Q$ let $dr_t=\kappa(\bar r-r_t)\,dt+\sigma\,dW_t$ with $\kappa>0$ and $\sigma>0$
\parencite{vasicek1977}. Then $P(t,T)=A(x)e^{-B(x)r_t}$ with $x=T-t$,
\[
  B(x)=\frac{1-e^{-\kappa x}}{\kappa},\qquad
  \ln A(x)=\bigl(B(x)-x\bigr)\Bigl(\bar r-\frac{\sigma^2}{2\kappa^2}\Bigr)-\frac{\sigma^2B(x)^2}{4\kappa}.
\]
\end{proposition}

\begin{proof}
The proof of the Vasicek short-rate proposition in Chapter~4 shows that, given $r_t$, $I=\int_t^Tr_u\,du$ is Gaussian with mean
$\bar rx+(r_t-\bar r)B(x)$ and variance $v=\sigma^2\int_0^xB(u)^2\,du$, where $B(u)$ is the function
written $\Lambda$ there. Hence $P(t,T)=\exp(-\E I+\tfrac12v)$, and integrating
$B(u)^2=(1-2e^{-\kappa u}+e^{-2\kappa u})/\kappa^2$ gives
$v=\tfrac{\sigma^2}{\kappa^2}\bigl(x-B(x)\bigr)-\tfrac{\sigma^2}{2\kappa}B(x)^2$. Collecting the terms
in $r_t$ and the rest gives the formula.
\end{proof}

The Gaussian law that makes this so easy also lets the rate go negative with positive probability.
The model of \textcite{cox1985} scales the noise by $\sqrt{r_t}$, so that the rate stays at or above
zero.

\begin{proposition}[CIR bond prices]\label{prop6:cir}
Under $Q$ let $dr_t=\kappa(\bar r-r_t)\,dt+\sigma\sqrt{r_t}\,dW_t$ with $\kappa,\bar r,\sigma>0$ and
$r_0\ge0$, and let $r$ be the nonnegative solution of this equation. With
$h=\sqrt{\kappa^2+2\sigma^2}$, define
\[
  B(x)=\frac{2(e^{hx}-1)}{2h+(\kappa+h)(e^{hx}-1)},\qquad
  A(x)=\Bigl(\frac{2h\,e^{(\kappa+h)x/2}}{2h+(\kappa+h)(e^{hx}-1)}\Bigr)^{2\kappa\bar r/\sigma^2}.
\]
Then $P(t,T)=A(T-t)e^{-B(T-t)r_t}$.
\end{proposition}

\begin{proof}
Direct differentiation shows that $A$ and $B$ solve the Riccati system
\[
  B'(x)=1-\kappa B(x)-\tfrac12\sigma^2B(x)^2,\qquad (\ln A)'(x)=-\kappa\bar r\,B(x),\qquad A(0)=1,\ B(0)=0,
\]
which the tests also check numerically. Both are nonnegative where needed: the numerator and denominator
of $B$ are positive for $x>0$, so $B\ge0$, and $\ln A$ starts at 0 and decreases, so $0<A\le1$.

Fix $T$, let $f(t,r)=A(T-t)e^{-B(T-t)r}$ and $M_t=e^{-\int_0^tr_u\,du}f(t,r_t)$ for $0\le t\le T$. With
$x=T-t$, $\partial_tf=f\,\bigl(-(\ln A)'(x)+B'(x)r\bigr)$, $\partial_rf=-Bf$ and $\partial_{rr}f=B^2f$, so
by It\^o's lemma the drift of $M$ is
\[
  M_t\Bigl[\bigl(-(\ln A)'-\kappa\bar rB\bigr)+r_t\bigl(B'+\kappa B+\tfrac12\sigma^2B^2-1\bigr)\Bigr]=0 ,
\]
and $dM_t=-B(T-t)\,\sigma\sqrt{r_t}\,M_t\,dW_t$. So $M$ is a local martingale\ix{Martingale}\ix{Martingale!local}. Because $r_t\ge0$, $B\ge0$ and
$A\le1$, it takes values in $[0,1]$, and a bounded local martingale is a martingale. Since
$M_T=e^{-\int_0^Tr_u\,du}$, \cref{eq6:bond} gives $P(t,T)=e^{\int_0^tr_u\,du}\,\E_Q[M_T\mid\mathcal F_t]=f(t,r_t)$.
\end{proof}

The equation has a unique nonnegative solution for these parameters, and when $r_0>0$ the solution
never reaches zero if $2\kappa\bar r\ge\sigma^2$ \parencite{feller1951,cox1985}. The proof needs only nonnegativity.

\subsection{Fitting the two models to the Canadian curve}

On a given date the market quotes bills for three months, six months and a year and benchmark bonds\ix{Benchmark bond}
for two to ten years, the eight yields of Chapter~4. Each model, with parameters
$(r_0,\kappa,\bar r)$ and a fixed $\sigma$, gives a discount function and so a model quote for each
instrument in its own convention: a simple yield $(1/P(0,t)-1)/t$ for a bill and a semi-annual par
yield $2\bigl(1-P(0,T)\bigr)/\sum_{i=1}^{2T}P(0,t_i)$ for a bond. The fit minimises the sum of squared
differences in basis points\ix{Basis point}. The curve's shape constrains $\kappa$ and $\bar r$, although how sharply it separates them depends on
the range of maturities and on the fitting criterion; it says little about $\sigma$, which enters only through a small convexity term, so $\sigma$ is set from history:
the standard deviation of the monthly changes in the 90 day bill yield between the \srVolN{} month-ends
in the file, times $\sqrt{12}$, which is \srSigV{} percentage points\ix{Percentage point}. For CIR\ix{Cox-Ingersoll-Ross model} the same size of move at the
average bill yield of \srMeanBill\% means $\sigma=\srSigC$. Using a real-world volatility under $Q$ is
an assumption; by Girsanov's theorem a change of measure\ix{Change of measure} moves the drift and leaves the volatility
alone, so it is a natural one.

\begin{table}[t]
\centering
\caption{\kind{historical} The Vasicek and CIR models fitted to eight Government of Canada bill and bond yields, with
$\sigma$ fixed from history. Errors are model minus market, in basis points. The parameters are
printed to five decimals, enough for the bond and swap values in the text to be reproduced from them.}
\label{tab6:fits}
\begin{tabular}{llrrrrr}
\toprule
date & model & $r_0$ (\%) & $\kappa$ & $\bar r$ (\%) & rms error & largest error\\
\midrule
\srDateNow & Vasicek & \srVasNowRx & \srVasNowKapx & \srVasNowRbarx & \srVasNowRms & \srVasNowMax\\
           & CIR     & \srCirNowRx & \srCirNowKapx & \srCirNowRbarx & \srCirNowRms & \srCirNowMax\\
\srDateInv & Vasicek & \srVasInvRx & \srVasInvKapx & \srVasInvRbarx & \srVasInvRms & \srVasInvMax\\
           & CIR     & \srCirInvRx & \srCirInvKapx & \srCirInvRbarx & \srCirInvRms & \srCirInvMax\\
\bottomrule
\end{tabular}
\end{table}

\Cref{tab6:fits,fig6:fits} show fits on two dates, the last date in the file and a date in July 2023
when the curve was inverted. On \srDateNow{} the fitted short rate\ix{Short rate} is \srVasNowR\% and its long-run level \srVasNowRbar\%, and both models fit the
eight quotes to \srVasNowRms{} basis points in root mean square. They
agree with each other to within \srGapNow{} basis points at every quote, because with this little
volatility the curve is almost entirely the deterministic pull of $r_0$ towards $\bar r$. The fitted
$\kappa$ of \srVasNowKap{} a year means a half-life of the pull, $\ln2/\kappa$, of \srHalfNow{} months. On \srDateInv{} the short rate sits far above its long-run level, and the inverted shape is
fitted less well, with errors up to \srVasInvMax{} basis points. One factor and three parameters can
make a curve that rises, falls or has one hump; they cannot follow every wiggle of a real curve, and
models used for pricing add a time-dependent drift so that today's curve is matched exactly. The
best known is the extended Vasicek\ix{Vasicek model} model of \textcite{hullwhite1990}, in which $\bar r$ becomes a
function of time chosen to reprice every zero-coupon bond on the day of the fit while bond prices\ix{Bond price} keep
the form $A e^{-Br}$ \parencite{brigo2006}.

\begin{figure}[t]
  \centering
  \includegraphics[width=0.94\textwidth]{ch06/fig04_curve_fits}
  \caption{\kind{historical} Government of Canada bill and benchmark yields (dots) and the model quotes of the fitted
  Vasicek and CIR models at the same maturities, on two dates. Bill quotes are simple yields and bond
  quotes semi-annual par yields, so each point compares like with like; the lines only join the points.}
  \label{fig6:fits}
\end{figure}

Where the two models part is in the distribution of future rates. Under the July 2023 Vasicek fit
the probability that the short rate is below zero in ten years is \srNegTen\%; in one year it is \srNegOneProb, zero for any practical purpose. CIR rates cannot go negative, and for the September 2026 fit the ratio
$2\kappa\bar r/\sigma^2$ is \srFeller, far above 1, the level below which the rate can reach zero. The case study avoids
choosing between the models by using the rates that actually occurred.

As a check on \cref{prop6:cir}, a simulation draws \srMcPaths{} paths of the September 2026 CIR model
from its exact noncentral chi-square transition and approximates $\int_0^{10}r_u\,du$ by the
trapezoid rule. With monthly steps (stream 2) the ten year bond comes out at \srMcMonth{} with
standard error\ix{Standard error} \srMcMonthSE, \srMcMonthZabs{} standard errors \srMcMonthDir{} the formula's \srMcExact; with weekly
steps (stream 3) it is \srMcWeek{} with standard error \srMcWeekSE, \srMcWeekZabs{} standard errors \srMcWeekDir. Neither
run shows a bias from the trapezoid rule that these standard errors can detect.

\subsection{Why a futures rate sits above the forward rate}

Chapter~4 showed that a futures\ix{Futures contract} price is the risk-neutral expectation of the underlying when
settlement is continuous, while a forward price\ix{Forward price} is an expectation under a different measure, and it
promised the size of the gap for contracts on rates. The cleanest case is the model of
\textcite{ho1986}, $dr_t=\eta(t)\,dt+\sigma\,dW_t$ under $Q$, with $\eta$ a deterministic function
chosen to fit today's curve.

\begin{proposition}[Convexity adjustment]\label{prop6:holee}
In the Ho-Lee\index{Ho-Lee} model, let $R(T_1,T_2)$ be the continuously compounded rate fixed at $T_1$ for the period
to $T_2$ and $f(0;T_1,T_2)$ the forward rate\ix{Forward rate} for that period,
\[
  R(T_1,T_2)=-\frac{\ln P(T_1,T_2)}{T_2-T_1},\qquad f(0;T_1,T_2)=\frac{\ln\bigl(P(0,T_1)/P(0,T_2)\bigr)}{T_2-T_1}.
\]
Then
\[
  \E_Q\bigl[R(T_1,T_2)\bigr]-f(0;T_1,T_2)=\tfrac12\sigma^2T_1T_2 .
\]
\end{proposition}

\begin{proof}
Write $\delta=T_2-T_1$ and $\Theta(t)=r_0+\int_0^t\eta(u)\,du$, so that $r_t=\Theta(t)+\sigma W_t$. As in
\cref{prop6:vasicek}, $\ln P(t,T)=-\E[I\mid\mathcal F_t]+\tfrac12\Var(I\mid\mathcal F_t)$ with
$I=\int_t^Tr_u\,du$. Here $\E[I\mid\mathcal F_t]=r_t(T-t)+\int_t^T(\Theta(u)-\Theta(t))\,du$ and
$\Var(I\mid\mathcal F_t)=\sigma^2(T-t)^3/3$. At $t=0$, $\ln P(0,T)=-\int_0^T\Theta(u)\,du+\sigma^2T^3/6$, so
$\delta f(0;T_1,T_2)=\int_{T_1}^{T_2}\Theta(u)\,du-\sigma^2(T_2^3-T_1^3)/6$. At $t=T_1$,
$\delta R(T_1,T_2)=\delta r_{T_1}+\int_{T_1}^{T_2}(\Theta(u)-\Theta(T_1))\,du-\sigma^2\delta^3/6$, whose
expectation, since $\E_Qr_{T_1}=\Theta(T_1)$, is $\int_{T_1}^{T_2}\Theta(u)\,du-\sigma^2\delta^3/6$.
Subtracting, the difference is $\sigma^2\bigl(T_2^3-T_1^3-\delta^3\bigr)/(6\delta)=\tfrac12\sigma^2T_1T_2$,
because $T_2^3-T_1^3-(T_2-T_1)^3=3T_1T_2(T_2-T_1)$.
\end{proof}

If a futures contract settled continuously on a price that is linear in $R(T_1,T_2)$, its rate would be
$\E_QR(T_1,T_2)$ by the futures and forward price proposition of Chapter~4, and the proposition would give the gap to the forward rate.
Real contracts, such as the three month CORRA\ix{CORRA} futures on the Montr\'eal Exchange, settle on a
compounded simple rate and are marked daily, so the formula is a guide to size. It grows with the
product of the two dates. With $\sigma=\srSigV$ points, a three month rate starting in one year
carries an adjustment of \hlAdjOne{} basis points; starting in two, five and ten years it is
\hlAdjTwo, \hlAdjFive{} and \hlAdjTen. Whether it matters depends on the horizon, the volatility and the precision needed: for near
contracts it is negligible, and at long horizons\ix{Horizon} it is material.


% =====================================================================================
\section{What do a swap and a cap cost on the fitted curve?}
\label{sec6:swaps}

An interest rate swap\ix{Swap rate} exchanges a fixed rate\ix{Fixed rate} for a floating one on a notional amount that never
changes hands. Take payment dates $T_0<t_1<\dots<t_n=T_n$ with accrual fractions
$\delta_i=t_i-t_{i-1}$ and $t_0=T_0$. At each $t_i$ the payer of the fixed rate $K$ pays $\delta_iK$
per unit of notional and receives $\delta_iL(t_{i-1},t_i)$, where $L(t_{i-1},t_i)=\bigl(1/P(t_{i-1},t_i)-1\bigr)/\delta_i$ is the
simple rate for the period, fixed at its start. Write
\[
  \mathcal A(T_0,T_n)=\sum_{i=1}^n\delta_iP(0,t_i)
\]
for the annuity\ix{Annuity} of the schedule, the value of receiving $\delta_i$ at each date.

\begin{proposition}[Swap rate]\label{prop6:swaprate}
Suppose the floating rate\ix{Floating rate} is the simple rate implied by the same discount curve\ix{Discount curve} that values the
payments. Then the floating leg is worth $P(0,T_0)-P(0,T_n)$, the payer swap is worth
$P(0,T_0)-P(0,T_n)-K\mathcal A(T_0,T_n)$, and the rate that makes it worth zero is
\[
  S_{\mathrm{sw}}(T_0,T_n)=\frac{P(0,T_0)-P(0,T_n)}{\mathcal A(T_0,T_n)} .
\]
\end{proposition}

\begin{proof}
The floating payment at $t_i$ is $1/P(t_{i-1},t_i)-1$. One unit of money held at $t_{i-1}$ buys
$1/P(t_{i-1},t_i)$ zero-coupon bonds\ix{Zero-coupon bond} maturing at $t_i$, so the first term is worth 1 at $t_{i-1}$ and
$P(0,t_{i-1})$ today; the second is worth $P(0,t_i)$. The sum over $i$ telescopes to
$P(0,T_0)-P(0,T_n)$. The fixed leg is $K\mathcal A$ by definition.
\end{proof}

The same telescoping works for a leg that compounds an overnight rate\ix{Overnight rate}, such as CORRA\ix{CORRA}, over each
period and pays at its end, provided payments are discounted at that overnight rate: rolling one unit
overnight from $t_{i-1}$ to $t_i$ produces exactly one plus the compounded payment. With $T_0=0$ and
semi-annual dates, $S_{\mathrm{sw}}$ is the par yield\ix{Par yield} $2\bigl(1-P(0,T)\bigr)/\sum_iP(0,t_i)$ of
\cref{sec6:short}, so a swap rate on a single curve is a par yield under another name. Market swap rates differ from government par yields, because swaps have their own floating rate
conventions and are discounted on a curve set by their collateral terms, while government bonds\ix{Sovereign bond} also
carry liquidity and other effects of the bond market; the numbers that follow use the fitted government curve as a stand-in.

On the September 2026 Vasicek fit of \cref{tab6:fits} the five year swap rate is \swFive\% and the
ten year rate \swTen\%. A swap that starts in five years and ends in ten has a forward rate\ix{Forward rate} of
\swFiveFive\% and an annuity of \swAnnFiveFive{} per unit of notional, so each basis point between
its fixed rate and the forward swap rate is worth \swAnnFiveFive{} basis points of notional today.
The annuity also measures the swap's rate risk. The payer swap is worth $\mathcal A\,(S_{\mathrm{sw}}-K)$. Holding the
annuity fixed, a small change in the forward swap rate changes the payer swap's value by the annuity
times that change; a general curve move also changes the annuity. The annuity plays the role that the
dollar duration $p(y)D$ played for a bond in \cref{sec6:duration}: it measures the change in value per
unit change in rate. A relative change would mean nothing here, since a swap struck at the forward rate
is worth zero.

\subsection{Caplets as options on a zero-coupon bond}

A caplet\ix{Interest rate cap} with strike\ix{Strike} $K$, fixing at $T$ and paying at $T+\delta$, pays
$\delta\bigl(L(T,T+\delta)-K\bigr)^+$ at $T+\delta$. A cap is a strip of caplets, and a borrower
with a floating rate loan who buys one has capped the reference rate in each period the cap covers,
though the loan's spread and any exposure the cap does not cover remain payable. The market quotes
caplets by the formula of \textcite{black1976}, with $F$ the forward rate\ix{Forward rate} for the period,
$\sigma_{\mathrm B}$ a volatility\ix{Volatility} and $d_1=\bigl(\ln(F/K)+\tfrac12\sigma_{\mathrm B}^2T\bigr)/(\sigma_{\mathrm B}\sqrt T)$:
\begin{equation}\label{eq6:black}
  \mathrm{Cpl}=\delta P(0,T+\delta)\bigl(F\Phi(d_1)-K\Phi(d_1-\sigma_{\mathrm B}\sqrt T)\bigr).
\end{equation}
The number $\sigma_{\mathrm B}$ that reproduces a price is its Black volatility\ix{Black volatility}. A model of the short rate gives the price directly, and for a Gaussian model\ix{Gaussian model} in closed form.

\begin{proposition}[Caplet as a bond put]\label{prop6:caplet}
\begin{enumerate}[label=(\roman*)]
  \item In any arbitrage-free model, the caplet is worth $1+\delta K$ European puts, expiring at $T$
  with strike $\bar P=1/(1+\delta K)$, on the zero-coupon bond\ix{Zero-coupon bond} maturing at $T'=T+\delta$.
  \item In the Vasicek\ix{Vasicek model} model of \cref{prop6:vasicek} with $\kappa>0$, such a put is worth
  \[
    \mathrm{ZBP}=\bar P\,P(0,T)\Phi(-h_P+\sigma_P)-P(0,T')\Phi(-h_P),\qquad
    h_P=\frac1{\sigma_P}\ln\frac{P(0,T')}{\bar P\,P(0,T)}+\frac{\sigma_P}2,
  \]
  where $\sigma_P=B(T'-T)\,\sigma\sqrt{(1-e^{-2\kappa T})/(2\kappa)}$.
\end{enumerate}
\end{proposition}

\begin{proof}
(i) The payment at $T'$ is known at $T$, so at $T$ it is worth
$P(T,T')\,\delta\bigl(L(T,T')-K\bigr)^+=\bigl(1-(1+\delta K)P(T,T')\bigr)^+=(1+\delta K)\bigl(\bar P-P(T,T')\bigr)^+$.

(ii) Let $Q^T$ be the measure with $dQ^T/dQ=e^{-\int_0^Tr_u\,du}/P(0,T)$. A payoff $Y$ at $T$ is
then worth $P(0,T)\E_{Q^T}Y$, and $\E_{Q^T}P(T,T')=P(0,T')/P(0,T)$. The $T$-maturity numeraire bond has
volatility coefficient $-B(T-t)\,\sigma$ under $Q$, which is deterministic, so by Girsanov's theorem $Q^T$ changes the drift of $r$ by a
deterministic function and leaves its volatility alone. Hence $r_T$ is Gaussian under $Q^T$ with the
same variance as under $Q$, $\sigma^2(1-e^{-2\kappa T})/(2\kappa)$, and
$\ln P(T,T')=\ln A(T'-T)-B(T'-T)r_T$ is Gaussian with variance $\sigma_P^2$. A lognormal variable with
mean $P(0,T')/P(0,T)$ and log variance $\sigma_P^2$ gives the put by the calculation behind Black's
formula\ix{Black's formula} with a zero rate, which is the expression stated.
\end{proof}

Part (ii) is the formula of \textcite{jamshidian1989}. For quarterly caplets struck at the forward on
the September 2026 fit, the forward rates starting in one, two and five years are \cpFOne\%,
\cpFTwo\% and \cpFFive\%, the caplets cost \cpAtmOne, \cpAtmTwo{} and \cpAtmFive{} basis points of
notional, and their Black volatilities are \cpVolOne\%, \cpVolTwo\% and \cpVolFive\%. The volatility
falls with expiry because the Vasicek rate is mean reverting: its standard deviation at $T$,
$\sigma\sqrt{(1-e^{-2\kappa T})/(2\kappa)}$, grows like $\sigma\sqrt T$ at first and then levels off
at $\sigma/\sqrt{2\kappa}$, and \cref{eq6:black} divides the spread of the rate by $\sqrt T$. It also
falls with strike. A Gaussian rate has the same spread in percentage points\ix{Percentage point} at every level, which in
proportion to the strike is large for a low strike and small for a high one; for the three year
caplet the Black volatility is \cpSkewLo\% at a 2\% strike and \cpSkewHi\% at 5\% (\cref{fig6:caplet}).
Market caplet volatilities show a skew or smile in strike whose direction and size vary with market
conditions, expiry and quoting convention, and the dependence of Black volatility on the level of
rates is one reason dealers often quote in
terms of normal volatility, the spread of the rate in basis points.

As a check on \cref{prop6:caplet}, a simulation draws \cpMcPaths{} values of $r_T$ and
$\int_0^Tr_u\,du$ from their exact joint Gaussian law under $Q$ (stream 4), prices the two year
caplet as the average of the discounted payoff\ix{Discounted payoff}, and gets \cpMc{} basis points with standard error
\cpMcSE, which is \cpMcZabs{} standard errors \cpMcDir{} the formula's \cpMcExact.

\begin{figure}[t]
  \centering
  \includegraphics[width=0.84\textwidth]{ch06/fig05_caplet_vols}
  \caption{\kind{exact} Black volatility of quarterly caplets priced in the September 2026 Vasicek fit, against
  the strike, for expiries of one, three and five years. The thin vertical lines mark the forward rate
  for each expiry.}
  \label{fig6:caplet}
\end{figure}

% =====================================================================================
\section{When is a firm's debt a short put on its assets?}
\label{sec6:merton}

The model of \textcite{merton1974} treats a firm as a pool of assets worth $V_t$ financed by equity
and by one zero-coupon debt issue with face value $K$ due at $T$. At $T$ the owners repay if
$V_T\ge K$ and hand over the assets otherwise. Equity is worth $(V_T-K)^+$ at $T$ and debt
$\min(V_T,K)=K-(K-V_T)^+$, so the lender holds a riskless bond and has written a put on the firm's
assets. Write the \emph{quasi-leverage}\index{Quasi-leverage} as $\ell=Ke^{-rT}/V_0$, the face of the debt discounted at the
riskless rate as a fraction of the assets.

\begin{proposition}[Merton spreads]\label{prop6:merton}
Suppose that under $Q$ the assets follow $dV_t=rV_t\,dt+\sigma_VV_t\,dW_t$ with constants $r$ and
$\sigma_V>0$, the firm pays nothing out before $T$, and default\ix{Default} can happen only at $T$. Let
$v=\sigma_V\sqrt T$, $d_1=(-\ln\ell+\tfrac12v^2)/v$ and $d_2=d_1-v$. Then equity is worth
$E_0=V_0\Phi(d_1)-Ke^{-rT}\Phi(d_2)$, the debt is worth
$Ke^{-rT}\bigl(\Phi(d_2)+\Phi(-d_1)/\ell\bigr)$, the risk-neutral probability of default is
$\Phi(-d_2)$, and the credit spread\ix{CDS spread}, the debt's continuously compounded yield less $r$, is
\[
  s(T)=-\frac1T\ln\Bigl(\Phi(d_2)+\frac{\Phi(-d_1)}{\ell}\Bigr),
\]
which depends on $\ell$, $\sigma_V$ and $T$ only.
\end{proposition}

\begin{proof}
Equity is a call on $V$ with strike $K$, so the formula of Chapter~5 with no dividend gives $E_0$,
since $\ln(V_0/K)+rT=-\ln\ell$. The debt is $V_0-E_0=V_0\Phi(-d_1)+Ke^{-rT}\Phi(d_2)$, and dividing by
$Ke^{-rT}$ gives the bracket. Default is $\{V_T<K\}$, whose $Q$ probability is $\Phi(-d_2)$ by the
lognormal law of $V_T$. The yield is $-\ln(\text{debt}/K)/T$, so the spread is the stated expression.
\end{proof}

\Cref{fig6:merton} draws $s(T)$ for four quasi-leverages with $\sigma_V=\mtSigV\%$, holding $\ell$
fixed as the maturity\ix{Maturity} changes. At $\ell=0.7$ the spread is \mtSSevenOne{} basis points at one year,
\mtSSevenFive{} at five and \mtSSevenTen{} at ten, with a peak of \mtPeakSeven{} near
\mtPeakTSeven{} years. At $\ell=0.9$ the curve is humped much earlier: \mtSNineOne{} at one year, a
peak of \mtPeakNine{} near \mtPeakTNine{} years, and \mtSNineTen{} at ten. Short maturities
behave very differently from long ones, and the limit can be stated exactly.

\begin{proposition}[Short maturities]\label{prop6:short}
Hold $\ell>0$ and $\sigma_V>0$ fixed and let $T\to0$.
\begin{enumerate}[label=(\roman*)]
  \item If $\ell<1$, then $T^{-j}s(T)\to0$ for every $j\ge0$.
  \item If $\ell=1$, then $\sqrt T\,s(T)\to\sigma_V/\sqrt{2\pi}$.
  \item If $\ell>1$, then $T\,s(T)\to\ln\ell$, so $s(T)\to\infty$.
\end{enumerate}
\end{proposition}

\begin{proof}
Let $u=1-\Phi(d_2)-\Phi(-d_1)/\ell=\Phi(-d_2)-\Phi(-d_1)/\ell$, the put divided by $Ke^{-rT}$, so that
$0\le u\le\Phi(-d_2)$ and $s(T)=-\ln(1-u)/T$. (i) Here $d_2=(-\ln\ell)/v-v/2\to\infty$. For $x\ge0$,
$\Phi(-x)\le\tfrac12e^{-x^2/2}$, so $u\le\tfrac12\exp\bigl(-(\ln\ell)^2/(2\sigma_V^2T)+O(1)\bigr)$, and
$-\ln(1-u)\le u/(1-u)$. The bound tends to zero faster than any power of $T$. (ii) Now $d_1=v/2$ and
$d_2=-v/2$, so $u=\Phi(v/2)-\Phi(-v/2)=v/\sqrt{2\pi}+O(v^3)$, and $-\ln(1-u)=u+O(u^2)$. Dividing by $T$
and multiplying by $\sqrt T$ gives the limit. (iii) Now $d_1,d_2\to-\infty$, so the bracket in
\cref{prop6:merton} tends to $1/\ell$ and $T\,s(T)\to\ln\ell$.
\end{proof}

A firm whose assets cover its discounted debt cannot default in a short time when the asset value\ix{Asset value}
moves continuously, so its short spreads vanish, and part (i) says they vanish very fast. At
$\ell=0.7$ the one month spread is \mtSSevenQ{} basis points. The limit arrives only once $v$ is small
beside $\ln(1/\ell)$, which for $\ell=0.9$ means maturities well under a month: its one month spread
is still \mtSNineQ. Observed short-dated corporate spreads stay positive where this model, with continuous asset paths and
default only at maturity, allows almost nothing, which is the
usual argument for adding jumps to the asset value or for modelling default directly, as the next
section does.

\begin{figure}[t]
  \centering
  \includegraphics[width=0.84\textwidth]{ch06/fig06_merton_spreads}
  \caption{\kind{exact} Credit spreads of \cref{prop6:merton} against maturity for quasi-leverage
  $\ell=0.5$, $0.7$, $0.9$ and $1.1$, with $\sigma_V=\mtSigV\%$ and $\ell$ held fixed along each curve.}
  \label{fig6:merton}
\end{figure}

\subsection{Backing out the assets from the equity}

Neither $V_0$ nor $\sigma_V$ is observed. The equity's value and volatility are, at least for a
listed firm, and the model ties them to the assets. Since $E_t=\mathcal C(t,V_t)$ for the call price $\mathcal C$,
It\^o's lemma gives $dE_t=\dots\,dt+\partial_V\mathcal C\,\sigma_VV_t\,dW_t$, and $\partial_V\mathcal C=\Phi(d_1)$, so the
instantaneous equity volatility\ix{Equity volatility} $\sigma_E$ satisfies
\begin{equation}\label{eq6:sige}
  \sigma_EE_0=\Phi(d_1)\,\sigma_VV_0 .
\end{equation}
In the model $\sigma_E$ changes as $V$ moves, so \cref{eq6:sige} uses the equity volatility at the
date of calibration. Together with $E_0=\mathcal C(0,V_0)$ it is two equations in $(V_0,\sigma_V)$.

Take a firm with equity worth 40, equity volatility 60\% and debt of face 80 due in five years, with
$r=\mtR\%$. Solving the two equations gives $V_0=\mcV$ and $\sigma_V=\mcSigV\%$, so the quasi-leverage
is \mcLev, the spread is \mcSpread{} basis points and the risk-neutral probability\ix{Probability} of default\ix{Default probability} is
\mcPdQ\%. The real-world probability needs the asset drift. With an expected asset return of
$\mu=\mcMu\%$, $\ln V_T$ has mean $\ln V_0+(\mu-\tfrac12\sigma_V^2)T$ under $\mathbb P$, the distance to
default $\bigl(\ln(V_0/K)+(\mu-\tfrac12\sigma_V^2)T\bigr)/(\sigma_V\sqrt T)$ is \mcDD{} standard
deviations, and the probability of default is \mcPdP\%. The gap between \mcPdQ\% and \mcPdP\% comes from the different asset drifts under $Q$ and $\mathbb P$;
within the model, and given the assumed $\mu$, it is the risk premium on the assets expressed as a
probability, and it is part of why a spread divided by the
loss given default overstates the default rate that lenders actually experience.

% =====================================================================================
\section{How is a default intensity read from credit default swaps?}
\label{sec6:hazard}

The structural model\ix{Merton model} explains default\ix{Default} by the assets. The reduced-form approach of
\textcite{jarrow1995}, \textcite{lando1998} and \textcite{duffie1999} gives up on the explanation and
describes the timing directly, which makes short spreads easy to match and fitting to market quotes
routine.

\begin{definition}\label{def6:hazard}
A default time $\tau$ has \emph{hazard rate\ix{Hazard rate}}\index{Hazard rate} $\lambda(t)\ge0$ under $Q$ if its survival function is
\[
  S(t)=Q(\tau>t)=e^{-H(t)},\qquad H(t)=\int_0^t\lambda(u)\,du .
\]
Then $Q(\tau\le t+h\mid\tau>t)=\lambda(t)h+o(h)$ at points where $\lambda$ is continuous, so $\lambda$ is
the rate of default per unit of time for a name that has survived to $t$.
\end{definition}

A credit default swap\ix{Credit default swap}\ix{Swap rate}\index{Credit default swap} on a name with maturity\ix{Maturity} $T$ has two legs. The protection buyer pays a spread
$s$ a year, quarterly, on the notional for as long as the name survives, plus the premium accrued
since the last payment date if it defaults. The seller pays the loss $1-R$ per unit of notional at
$\tau$ if $\tau\le T$, where $R$ is the recovery rate\ix{Recovery rate}. If $\tau$ is independent of interest rates\ix{Interest rate} under
$Q$, so that a payment at $u$ contingent on survival is worth $P(0,u)S(u)$, the two legs are worth
$s\,\mathcal A_\tau(T)$ and $\mathrm{Prot}(T)$, where
\begin{align*}
  \mathcal A_\tau(T)&=\sum_{i=1}^{4T}\tfrac14P(0,t_i)S(t_i)+\sum_{i=1}^{4T}\int_{t_{i-1}}^{t_i}(u-t_{i-1})P(0,u)\lambda(u)S(u)\,du,\\
  \mathrm{Prot}(T)&=(1-R)\int_0^TP(0,u)\lambda(u)S(u)\,du ,
\end{align*}
with $t_i=i/4$. The par spread\ix{Par spread} is $\mathrm{Prot}(T)/\mathcal A_\tau(T)$. The code evaluates each integral
by the midpoint rule with 24 steps a quarter.

\begin{proposition}[Credit triangle]\label{prop6:triangle}
If $\lambda$ is constant and the premium is paid continuously, so that the premium leg is
$s\int_0^TP(0,u)S(u)\,du$, then the par spread is $s=\lambda(1-R)$ for every maturity $T$ and every
discount curve\ix{Discount curve}.
\end{proposition}

\begin{proof}
With constant $\lambda$, $\mathrm{Prot}(T)=(1-R)\lambda\int_0^TP(0,u)S(u)\,du$, which is $(1-R)\lambda$ times
the continuous premium leg per unit of spread.
\end{proof}

With quarterly payments and accrued premium the identity is close: on the September 2026 curve a flat
hazard of 1\% and recovery of 40\% give a five year par spread of \hzFlatPar{} basis points, against
\hzFlatTri{} from the triangle.

Market quotes come at several maturities, and a piecewise-constant hazard with one level per
interval fits them exactly. Order the quoted maturities $T_1<\dots<T_m$ and let $\lambda$ equal $\lambda_j$
on $(T_{j-1},T_j]$, with the last level extended beyond $T_m$. The $T_1$ quote involves only $\lambda_1$;
given $\lambda_1$, the $T_2$ quote involves only $\lambda_2$; and so on, so the curve is found one
level at a time with a scalar root finder. \Cref{tab6:hazard} applies this to an upward sloping curve
of quotes of \hzQOne, \hzQThree, \hzQFive, \hzQSeven{} and \hzQTen{} basis points at one, three, five,
seven and ten years, with $R=\hzRec\%$, discounting on the September 2026 Vasicek curve. The fitted
curve\ix{Curve fitting} reprices all five quotes, which a test checks.

\begin{table}[t]
\centering
\caption{\kind{exact} Hazard rates bootstrapped from a CDS curve with recovery \hzRec\%. The last column is the
level of the hazard on the interval that ends at the maturity in the first column; the level for
seven to ten years is used beyond ten.}
\label{tab6:hazard}
\begin{tabular}{rrrr}
\toprule
maturity (years) & quote (bp) & $s/(1-R)$ (\%) & hazard on the interval (\%)\\
\midrule
1  & \hzQOne   & \hzTriOne   & \hzLOne\\
3  & \hzQThree & \hzTriThree & \hzLThree\\
5  & \hzQFive  & \hzTriFive  & \hzLFive\\
7  & \hzQSeven & \hzTriSeven & \hzLSeven\\
10 & \hzQTen   & \hzTriTen   & \hzLTen\\
\bottomrule
\end{tabular}
\end{table}

The triangle values in the third column are averages over the whole life of each contract, and the
bootstrapped levels are forward values for each interval, so on a rising curve the levels climb above
the averages: the five year quote implies \hzTriFive\% by the triangle, the hazard between three and
five years is \hzLFive\%, and the average hazard over five years, $H(5)/5$, is \hzAvgFive\%.
\Cref{fig6:hazard} shows the step function and the survival curve. The five year survival
probability is \hzSFive\%, so the risk-neutral probability of default within five years is
\hzPdFiveX\%.

\begin{figure}[t]
  \centering
  \includegraphics[width=0.84\textwidth]{ch06/fig07_hazard}
  \caption{\kind{exact} The bootstrapped hazard rate of \cref{tab6:hazard} (steps), the triangle values
  $s_T/(1-R)$ at each quoted maturity (dots), and the survival probability $S(t)$ on the right axis
  (dashed).}
  \label{fig6:hazard}
\end{figure}

Default times with a given hazard are easy to simulate.

\begin{proposition}[Simulating a default time]\label{prop6:deftime}
Let $H$ be continuous and nondecreasing with $H(0)=0$, let $\xi$ be a standard exponential variable,
and let $\tau=\inf\{t\ge0:H(t)\ge\xi\}$, with $\tau=\infty$ if the set is empty. Then
$Q(\tau>t)=e^{-H(t)}$ for every $t\ge0$.
\end{proposition}

\begin{proof}
Because $H$ is continuous and nondecreasing, $\tau>t$ holds exactly when $H(t)<\xi$, and
$Q(\xi>H(t))=e^{-H(t)}$.
\end{proof}

For a piecewise-constant hazard the infimum is found interval by interval. With \hzMcN{} draws
(stream 5), the share of defaults within five years is \hzMcPd\% with standard error \hzMcPdSE,
\hzMcPdZ{} standard errors from \hzPdFiveX\%. The average of $(1-R)P(0,\tau)\mathbf 1\{\tau\le5\}$ is
\hzMcProt{} basis points of notional with standard error \hzMcProtSE, against \hzProt{} from the
formula for $\mathrm{Prot}(5)$, a gap of \hzMcProtZ{} standard errors. The independence of $\tau$ and
the rates is what lets the formula multiply $P(0,u)$ by $S(u)$; a name whose default risk rises when
rates fall would need a joint model.

\subsection{How single-name contracts are quoted}
\label{sec6:cdsquote}

The par spread\ix{Par spread} is the natural output of the model, but since 2009 North American single-name contracts
have not traded at par. They carry a fixed running coupon\ix{Coupon} of 100 basis points\ix{Basis point}, usual for investment
grade names, or 500 for high yield, and the difference is settled by an upfront payment
\parencite{markit2009}\index{Credit default swap!upfront payment}. The buyer also pays a full coupon on each quarterly date,
and the seller rebates at the trade the premium accrued before it. With a coupon $C$ the upfront paid by
the buyer per unit of notional is
\[
  U=\mathrm{Prot}(T)-C\,\mathcal A_\tau(T)=\bigl(s-C\bigr)\mathcal A_\tau(T),
\]
where $s$ is the par spread on the same hazard curve, so a name trading inside its coupon has a negative
upfront and the buyer receives cash. Dealers quote a 100 basis point name by a conventional spread, and
the upfront is computed from it with a standard model that assumes one flat hazard\ix{Flat hazard} for all horizons and
standard recoveries; a 500 basis point name is usually quoted in points upfront.

For the five year quote of \hzQFive{} basis points the flat hazard that reprices it is \cdFlatLam\%, the
risky annuity\ix{Annuity} is \cdAnnFlat{} years, and on a 100 basis point coupon the buyer receives \cdUpIg{} per
100 of notional. Valued on the bootstrapped curve\ix{Bootstrapping a curve} of \cref{tab6:hazard} instead, the annuity is
\cdAnnCurve{} and the buyer receives \cdUpIgCurve, a gap of \cdUpIgGapBp{} basis points of notional
that comes from the flat hazard putting more of the default risk early. A high yield name quoted at
\cdQHy{} basis points has a flat hazard of \cdLamHy\%, an annuity of \cdAnnHy{} and, on the 500 basis
point coupon, an upfront of \cdUpHy{} points paid by the buyer; on a 100 basis point coupon it would
need \cdUpHyOne. The code checks the identity above on both curves and that the upfront is zero
when the quoted spread equals the coupon. Real conversions also fix the discount curve\ix{Discount curve}, the trade and
settlement dates and the day count, which the sketch here leaves out.

% =====================================================================================
\section{Why do defaults in a pool arrive together?}
\label{sec6:copula}

A loan fund, a CLO\ix{Collateralised loan obligation} or a bank's book holds many names. If their defaults\ix{Default} were independent\ix{Independence}, the share
that defaults in a large pool\ix{Large homogeneous pool} would be close to its expectation, so a big pool would carry little risk beyond its expected loss.
Real pools carry more, because borrowers share an economy. The standard way to put that into numbers is a
one-factor\ix{Common factor} Gaussian model\ix{Gaussian model}, the form of the Gaussian copula\ix{Copula}\index{Gaussian copula} of \textcite{li2000} that the market adopted.

\begin{definition}\label{def6:copula}
In a pool of $n$ names, name $i$ defaults by the horizon if
$X_i=\sqrt\rho\,M+\sqrt{1-\rho}\,\varepsilon_i\le\Phi^{-1}(p)$, where $M,\varepsilon_1,\dots,\varepsilon_n$
are independent standard normal variables and $0\le\rho<1$. Each name then defaults with
probability\ix{Probability} $p$, and any two latent variables $X_i,X_j$ have correlation $\rho$. Each default loses a
fraction $1-R$ of an equal share of the pool, so the pool's loss is $L=(1-R)N/n$ with $N$ the number
of defaults.
\end{definition}

\begin{proposition}[Conditional independence]\label{prop6:cond}
Given $M=m$, defaults are independent with probability
\[
  p(m)=\Phi\Bigl(\frac{\Phi^{-1}(p)-\sqrt\rho\,m}{\sqrt{1-\rho}}\Bigr),
\]
and $Q(N=k)=\int_{-\infty}^{\infty}\binom nk p(m)^k\bigl(1-p(m)\bigr)^{n-k}\varphi(m)\,dm$ for
$k=0,\dots,n$, where $\varphi$ is the standard normal density.
\end{proposition}

\begin{proof}
Given $M=m$, $X_i\le\Phi^{-1}(p)$ exactly when $\varepsilon_i\le(\Phi^{-1}(p)-\sqrt\rho\,m)/\sqrt{1-\rho}$,
and the $\varepsilon_i$ are independent of each other and of $M$. So $N$ is binomial given $M$, and
averaging over $M$ gives the formula.
\end{proof}

Take a pool of \cpN{} names, each with the five year default probability of \cref{sec6:hazard},
\hzPdFiveX\% (\hzPdFive\% to two places; every calculation below uses the unrounded value), and recovery \hzRec\%, so the expected pool loss is \cpEL\% at every $\rho$.
\Cref{fig6:counts} shows how the distribution of the number of defaults changes with $\rho$. With
independent names the probability of no default is \cpPZeroZero\% and the probability of 16 or more
is \cpPTailZero\%. At $\rho=0.1$ these are \cpPZeroOne\% and \cpPTailOne\%; at $0.3$,
\cpPZeroThree\% and \cpPTailThree\%; and at $0.6$, \cpPZeroSix\% and \cpPTailSix\%. The mean stays
where it was while the mass moves to the two ends: correlation makes a quiet five years and a
disastrous one both more likely.

\begin{figure}[t]
  \centering
  \includegraphics[width=0.84\textwidth]{ch06/fig08_default_counts}
  \caption{\kind{exact} Probability of each number of defaults within five years among \cpN{} names with default
  probability \hzPdFiveX\%, for factor correlations $\rho=0$, $0.1$, $0.3$ and $0.6$, by
  \cref{prop6:cond}. Counts above 40 are not drawn; at $\rho=0.6$ they carry \cpPAboveSix\% of the probability.}
  \label{fig6:counts}
\end{figure}

As the pool grows the binomial noise averages away and only the factor is left. The limit is due to
\textcite{vasicek2002}.

\begin{proposition}[Large pool]\label{prop6:lhp}
Let $0<p<1$ and $0<\rho<1$. As $n\to\infty$ the default share $N/n$ converges\ix{Convergence} almost surely to $p(M)$,
and for $0<x<1$
\[
  Q\bigl(p(M)\le x\bigr)=\Phi\Bigl(\frac{\sqrt{1-\rho}\,\Phi^{-1}(x)-\Phi^{-1}(p)}{\sqrt\rho}\Bigr).
\]
\end{proposition}

\begin{proof}
Given $M=m$ the default indicators are independent Bernoulli variables with mean $p(m)$, so by the
strong law of large numbers\ix{Law of large numbers} $N/n\to p(m)$ with conditional probability one. Integrating the
conditional probability of convergence, which is 1 for every $m$, over the law of $M$ gives
convergence with probability one. Since $p(m)$ is decreasing in $m$, $p(M)\le x$ exactly when
$M\ge\bigl(\Phi^{-1}(p)-\sqrt{1-\rho}\,\Phi^{-1}(x)\bigr)/\sqrt\rho$, whose probability is the expression
stated.
\end{proof}

\subsection{Cutting the pool into tranches}

A tranche\ix{Tranche}\index{Tranche} with attachment $a$ and detachment $d$, $0\le a<d\le1$, absorbs the part of the pool loss
between $a$ and $d$. Its loss as a fraction of its size is
\[
  L_{[a,d]}=\frac{\min(L,d)-\min(L,a)}{d-a}=\frac{\min\bigl((L-a)^+,\,d-a\bigr)}{d-a}.
\]
The \emph{equity} tranche starts at $a=0$ and takes the first losses; a \emph{senior} tranche ends at
$d=1$ and takes the last.

\begin{proposition}[What tranching preserves and what it moves]\label{prop6:tranche}
\begin{enumerate}[label=(\roman*)]
  \item For any law of $L$ and any attachment points $0=a_1<d_1=a_2<\dots<d_m=1$,
  \[\textstyle\sum_j(d_j-a_j)\E L_{[a_j,d_j]}=\E L.\]
  \item Let $L$ and $L'$ take values in $[0,1]$ with $\E f(L')\ge\E f(L)$ for every convex $f$. Then
  the equity tranche $[0,d]$ has an expected loss\ix{Expected loss} under $L'$ no larger than under $L$, and the senior
  tranche\ix{Senior debt} $[a,1]$ has one no smaller.
\end{enumerate}
\end{proposition}

\begin{proof}
(i) The sum of $\min(L,d_j)-\min(L,a_j)$ telescopes to $\min(L,1)-\min(L,0)=L$. (ii) The equity loss
$\min(L,d)$ is concave in $L$, so $-\min(L,d)$ is convex, and the senior loss $(L-a)^+$ is convex.
\end{proof}

The condition in (ii) says that $L'$ is more spread out than $L$ in the convex order; it forces equal
means, by taking $f(x)=x$ and $f(x)=-x$. Part (i) says tranching only redistributes the pool's expected
loss. Part (ii) says that anything spreading the pool loss out moves expected loss from the equity to
the senior notes, and in every case computed here a higher $\rho$ acts that way. \Cref{tab6:tranches}
gives the expected losses of six tranches at three correlations for the pool of \cpN{} names and in
the large-pool limit, computed by quadrature\ix{Numerical integration} over $M$; the thickness-weighted sum of each column is
\cpEL\%, as (i) requires, and a test checks it. As $\rho$ goes from 0.1 to 0.6 the equity loss falls
from \trEqOne\% to \trEqSix\% and the 30 to 100\% tranche goes from \trSupOne\% to \trSupSix\%. The
mezzanine tranches have payoffs that are neither convex nor concave, so (ii) says nothing about them:
the 3 to 7\% tranche falls from \trMzAOne\% to \trMzASix\% and the 7 to 10\% tranche rises from
\trMzBOne\% to \trMzBSix\%. The two pools differ
most at low correlation, where the 100-name pool still has binomial noise: \trEqOne\% against
\lhEqOne\% for the equity.

\begin{table}[t]
\centering
\caption{\kind{exact} Five year expected loss of each tranche, as a percentage of the tranche, for \cpN{} names
with default probability \hzPdFiveX\%, $1-S(5)$ from \cref{tab6:hazard}, and recovery \hzRec\% (pool) and in the large-pool limit of
\cref{prop6:lhp} (limit). An entry of 0.0000 is below 0.00005\%.}
\label{tab6:tranches}
\begin{tabular}{lrrrrrr}
\toprule
 & \multicolumn{2}{c}{$\rho=0.1$} & \multicolumn{2}{c}{$\rho=0.3$} & \multicolumn{2}{c}{$\rho=0.6$}\\
\cmidrule(lr){2-3}\cmidrule(lr){4-5}\cmidrule(lr){6-7}
tranche (\%) & pool & limit & pool & limit & pool & limit\\
\midrule
0 to 3    & \trEqOne  & \lhEqOne  & \trEqThree  & \lhEqThree  & \trEqSix  & \lhEqSix\\
3 to 7    & \trMzAOne & \lhMzAOne & \trMzAThree & \lhMzAThree & \trMzASix & \lhMzASix\\
7 to 10   & \trMzBOne & \lhMzBOne & \trMzBThree & \lhMzBThree & \trMzBSix & \lhMzBSix\\
10 to 15  & \trSnAOne & \lhSnAOne & \trSnAThree & \lhSnAThree & \trSnASix & \lhSnASix\\
15 to 30  & \trSnBOne & \lhSnBOne & \trSnBThree & \lhSnBThree & \trSnBSix & \lhSnBSix\\
30 to 100 & \trSupOne & \lhSupOne & \trSupThree & \lhSupThree & \trSupSix & \lhSupSix\\
\bottomrule
\end{tabular}
\end{table}

\begin{figure}[t]
  \centering
  \includegraphics[width=0.84\textwidth]{ch06/fig09_tranche_el}
  \caption{\kind{exact} Expected loss of the six tranches of \cref{tab6:tranches} for the \cpN-name pool against
  the factor correlation, on a log scale. Values below 0.0001\% are off the scale.}
  \label{fig6:tranche}
\end{figure}

The senior tranches are hit only in the far tail of $M$, and a quadrature that does not know where the
kinks of the tranche payoff are can step over the region entirely and return zero. The code splits
the range of $M$ at the two points where the conditional loss $(1-R)p(m)$ equals $a$ and $d$, and that split calculation is the one behind the tranche results above. The
listing below is a simpler cross-check, a self-contained version of the large-pool calculation for
$\rho=0.3$ with $p=0.05753$, the default probability to five places.
\Needspace*{18\baselineskip}
\lstinputlisting[language=Python]{code/ch06/snippets/lhp_tranche.py}
The listing integrates over a fixed range without splitting it. For these six tranches it agrees with the split calculation, which a test checks, but it is not a
general method: for thinner or more senior tranches only the split is reliable.

As a check on the 100-name column at $\rho=0.3$, a simulation draws \cpMcScen{} scenarios of $M$ and
the $\varepsilon_i$ (stream 6). The equity tranche comes out at \cpMcEq\% with standard error\ix{Standard error}
\cpMcEqSE{} and the 15 to 30\% tranche at \cpMcSnB\% with standard error \cpMcSnBSE. The largest of
the six gaps, \cpMcMaxZ{} standard errors, is in the \cpMcMaxName{} tranche, \cpMcMaxEst\% with standard
error \cpMcMaxSE{} against \cpMcMaxTab\% in the table, and \cpMcSignNeg{} of the six are negative. The six
estimates are far from independent, since they share every draw of $M$, so a set of factor draws that
happens to be mild pulls all of them down together. A second run on stream 8 has a largest gap of
\cpMcMaxZAlt{} standard errors.

\subsection{How a structure is built, and why its maturities must line up}\label{sec6:structures}

A cash collateralised loan obligation\ix{Collateralised loan obligation} buys a portfolio of loans\ix{Loan}, often a hundred or more, and pays for
it by issuing notes in order of seniority with a thin layer of equity at the bottom. Interest from
the loans pays the notes' coupons\ix{Coupon} from the top down, and principal pays them down in the same order.
Coverage tests protect the senior notes: if the par value of the loans falls too far relative to the
notes at and above a given class, or if interest received falls too far below interest due, cash that
would have gone to the junior notes is diverted to repay the senior ones. A synthetic structure
holds no loans: it takes credit risk\ix{Credit risk} on a reference pool through credit default swaps, and its
investors in effect sell protection on tranches\ix{Tranche} of that pool, and for a homogeneous pool like the one
above the one-factor\ix{Common factor} model gives those tranches the expected losses of \cref{tab6:tranches}.

The tranche losses above are losses by a five year horizon. A note is exposed until it is repaid, and
when it is repaid depends on the timing of the collateral's cash. This is where maturity\ix{Maturity} matching
enters. If the notes have a legal final maturity before the collateral matures, the vehicle must sell
or refinance\ix{Refinancing risk} loans to repay them, at whatever price the market offers on that date. A legal final maturity later than the collateral's reduces that risk, because scheduled principal from
the loans can repay the notes as it arrives; defaults, extensions, the timing of prepayments\ix{Prepayment},
reinvestment\ix{Reinvestment} limits and the waterfall can still leave a shortfall or force a sale. Cash flow
structures therefore usually set the legal final maturity of the notes after the latest maturity
the collateral is allowed to have, and limit what the manager may buy during the reinvestment period to loans maturing before that date.
The rate side is matched separately: loans and notes in a typical loan obligation both pay floating
coupons, so a move in the general level of rates changes both legs and the net duration of \cref{sec6:duration} is usually small, although fixed rate
loans, rate floors and reset lags leave some exposure. What remains is the risk of wider credit spreads\ix{Credit spread} and the timing of the collateral's principal
payments.

The funding problems of asset-backed commercial paper conduits and structured investment vehicles in
2007 and 2008 show the danger. They held assets with lives of years and funded them with paper of a
few months, so they had to roll that funding constantly; when buyers stopped buying, many had to sell
into a falling market. Canada had its own version: in August 2007 the market for third-party
asset-backed commercial paper froze, and holders waited until the Montreal Accord restructuring
turned the paper into notes maturing years later. Whether a borrower can refinance when its funding
comes due is the same question, at the level of a single loan, that the case study in
\cref{sec6:case} asks of commercial mortgages\ix{Mortgage}\ix{Commercial mortgage}.

% =====================================================================================
\section{What does a private lender's screen measure?}
\label{sec6:private}

Private credit\ix{Private credit} funds lend to mid-sized companies and against commercial property, and before any
model is run a deal is usually screened on a few ratios. This section takes the ratios used by the
open-source private credit screener \texttt{private-credit-screener} \parencite{pcscreener}, a small browser tool that scores a loan
on coverage, leverage and collateral and grades it, and asks what they measure. The code here is a
Python port of the screener's calculation file, and the tests run the original JavaScript on the same
inputs to confirm that the two give the same scores and grades, and the same metrics up to floating-point
rounding.

\begin{definition}\label{def6:screen}
A loan of amount $\mathcal L$ against a property of value $V$ with net operating income\ix{Net operating income} $\mathcal N$ a
year carries an annual rate $r_{\mathcal L}$, compounded monthly, and amortises over $n_{\mathrm{am}}$ years or pays
interest only\ix{Interest only}. Its annual debt service\ix{Debt service} is $\mathcal L r_{\mathcal L}$ if interest only, and otherwise
$12\mathcal L r_{\mathrm m}/\bigl(1-(1+r_{\mathrm m})^{-12n_{\mathrm{am}}}\bigr)$ with $r_{\mathrm m}=r_{\mathcal L}/12$; dividing by $\mathcal L$ gives the
\emph{loan constant\ix{Loan constant}}\index{Loan constant} $k$. The screen computes
\[
  \mathrm{LTV}=\frac{\mathcal L}{V},\qquad \mathrm{DSCR}=\frac{\mathcal N}{k\mathcal L},\qquad
  \mathrm{DY}=\frac{\mathcal N}{\mathcal L},\qquad
  \text{exit LTV}=\frac{b}{\mathcal N(1+g)^T/(q_{\mathrm{in}}+\Delta q)},
\]
the loan to value\ix{Loan to value}\index{Loan to value}, debt service coverage\index{Debt service coverage ratio}, debt yield\ix{Debt yield} and loan to value at maturity. Here $b$ is the
balance due at the loan's term of $T$ years, $g$ is the annual growth of income, compounded annually,
$q_{\mathrm{in}}$ is the going-in capitalisation rate\ix{Capitalisation rate} entered for the property and $\Delta q$ the
assumed rise in that rate by the exit.
\end{definition}

Three of the four ratios are less independent than they look.

\begin{proposition}[Two identities]\label{prop6:identities}
Let $q=\mathcal N/V$ be the cap rate implied by the income and the value. Then
$\mathrm{LTV}\times\mathrm{DY}=q$ and $\mathrm{DSCR}=\mathrm{DY}/k$.
\end{proposition}

\begin{proof}
$(\mathcal L/V)(\mathcal N/\mathcal L)=\mathcal N/V$, and $\mathcal N/(k\mathcal L)=(\mathcal N/\mathcal L)/k$.
\end{proof}

Given the cap rate and the loan constant, the loan to value fixes the other two, so the first three ratios are not independent measures of credit quality: through the cap rate and the
loan constant they reduce to how much was lent, how the market prices\ix{Market price} the income, and how expensive
the debt is. The screen maps each ratio to a score between 0 and 1 along a straight line
that is clipped at the ends: loan to value from 85\% down to 50\%, coverage from 1 up to 1.6, debt
yield from 6\% up to 14\%, exit loan to value from 90\% down to 50\%, and a sponsor rating from 1 to
5. It weights them 0.25, 0.25, 0.15, 0.25 and 0.10, rounds the total to a whole number out of 100,
and grades 75 and above A, 60 B, 45 C and the rest D.

\Cref{tab6:screen} applies the screen to the three made-up example loans that ship with the tool: a
stabilised industrial building in Calgary, a condominium inventory loan in Toronto and a mixed-use
bridge loan in Ottawa. The Calgary loan is interest only at 9.5\%, so its constant equals its rate and
its coverage is its debt yield divided by \pcCalConst\%. The Toronto loan's income covers just over half its debt service, which the coverage of \pcTorDscr{} shows and the debt yield of \pcTorDy\% shows
again. The Ottawa loan has a healthier debt yield, \pcOttDy\%, but amortises over 25 years at
10.25\%, so its constant is \pcOttConst\% and its coverage falls below 1. The exit test uses the cap
rate entered for the property, which for Toronto and Ottawa is not the one the income and value
imply.

\begin{table}[t]
\centering
\caption{\kind{exact} The screener's metrics, score and grade for its three example loans. Ratios are in percent
except the coverage; the cap rate is $\mathcal N/V$ and the constant is $k$.}
\label{tab6:screen}
\begin{tabular}{lrrrrrrrl}
\toprule
loan & LTV & DSCR & DY & cap rate & constant & exit LTV & score & grade\\
\midrule
Calgary & \pcCalLtv & \pcCalDscr & \pcCalDy & \pcCalCap & \pcCalConst & \pcCalExitLtv & \pcCalScore & \pcCalGrade\\
Toronto & \pcTorLtv & \pcTorDscr & \pcTorDy & \pcTorCap & \pcTorConst & \pcTorExitLtv & \pcTorScore & \pcTorGrade\\
Ottawa  & \pcOttLtv & \pcOttDscr & \pcOttDy & \pcOttCap & \pcOttConst & \pcOttExitLtv & \pcOttScore & \pcOttGrade\\
\bottomrule
\end{tabular}
\end{table}

\subsection{From a ratio to an expected loss}

A screen ranks loans, but it says nothing about how likely a loss is or how large it would be.
For that one needs a law for the collateral. A lender holding the slice of a property's value from
$b_{\mathrm{lo}}$ to $b_{\mathrm{hi}}$, such as a mezzanine\ix{Subordination} lender behind a senior loan, recovers from a sale at
$Y=(1-\theta)V_T$ after selling costs\ix{Selling costs} $\theta$, and the senior lender is the case $b_{\mathrm{lo}}=0$.

\begin{proposition}[Slice expected loss]\label{prop6:slice}
A claim with face $b_{\mathrm{hi}}-b_{\mathrm{lo}}$ that is paid from $Y$ after claims of $b_{\mathrm{lo}}$ loses
$(b_{\mathrm{hi}}-Y)^+-(b_{\mathrm{lo}}-Y)^+$. If $Y$ is lognormal with mean $m$ and log standard deviation
$v$, its expected loss\ix{Expected loss} is $\Pi(b_{\mathrm{hi}})-\Pi(b_{\mathrm{lo}})$, where $\Pi(0)=0$ and, for $K>0$,
\[
  \Pi(K)=\E(K-Y)^+=K\Phi(-d_2)-m\Phi(-d_1),\qquad d_1=\frac{\ln(m/K)+\tfrac12v^2}{v},\quad d_2=d_1-v .
\]
\end{proposition}

\begin{proof}
The claim receives $\min\bigl((Y-b_{\mathrm{lo}})^+,\,b_{\mathrm{hi}}-b_{\mathrm{lo}}\bigr)$, and its face less this
is $(b_{\mathrm{hi}}-Y)^+-(b_{\mathrm{lo}}-Y)^+$, as checking the three cases $Y\le b_{\mathrm{lo}}$,
$b_{\mathrm{lo}}<Y<b_{\mathrm{hi}}$ and $Y\ge b_{\mathrm{hi}}$ shows. The expectation of each term is the put
of Chapter~5 with a zero rate and forward $m$.
\end{proof}

The slice uses the attachment and detachment logic of the tranches\ix{Tranche} in \cref{sec6:copula}, and its loss
is set by what the property fetches on sale in place of the share of a pool that is lost. Take a property worth 100 whose value at $T=\slT$ years has mean $100(1+\slG\%)^{\slT}$
and log standard deviation $\sigma_V\sqrt T$, selling costs of \slCost\%, and loans at \slRate\% that
amortise over \slAmort{} years. A senior loan at 60\% loan to value\ix{Loan to value} has a balance of \slBs{} at five
years, and a loan to 75\% has \slBm, so the mezzanine slice runs from \slBs{} to \slBm. With
$\sigma_V=10\%$ the five year expected losses are \slSenTen\% of face for the senior loan and
\slMezTen\% for the mezzanine slice; with 15\% they are \slSenFifteen\% and \slMezFifteen\%; with
20\%, \slSenTwenty\% and \slMezTwenty\%. The probability that the sale does not cover the senior
balance goes from \slPdSenTen\% to \slPdSenTwenty\% over the same range, and for the mezzanine
balance from \slPdMezTen\% to \slPdMezTwenty\%. Doubling the volatility of the collateral multiplies the senior loss by \slSenRatio{} and the
mezzanine loss by \slMezRatio, while every ratio in
\cref{def6:screen} stays where it was. \Cref{fig6:slice} draws the losses against the loan to value
at which each claim stops.

\begin{figure}[t]
  \centering
  \includegraphics[width=\textwidth]{ch06/fig10_slice_el}
  \caption{\kind{exact} Five year expected loss by \cref{prop6:slice} for a property worth 100 whose expected value grows
  \slG\% a year, selling costs of \slCost\% and loans at \slRate\% amortising over \slAmort{}
  years, for three values of $\sigma_V$. Left: a senior loan from 0 up to the loan to value on the
  axis. Right: a slice 15 points thick ending at the loan to value on the axis.}
  \label{fig6:slice}
\end{figure}

% =====================================================================================
\section{How does a direct lender earn its spread, and what do its returns hide?}
\label{sec6:direct}

The loans of \cref{sec6:private} are secured on property. Most private credit\ix{Private credit} is lent to companies.
\textcite{cai2024} describe it as debt that is not publicly traded, provided by non-bank entities such
as private credit funds and business development companies to private businesses, usually middle-market
firms with annual revenues between \$10 million and \$1 billion. A loan negotiated between one borrower
and one lender, or a small group of lenders, is called direct lending\ix{Direct lending}. Such loans are usually senior
secured\ix{Senior secured loan}, almost all pay a floating rate, and they are held to maturity or to a refinancing because there
is no liquid secondary market. The International Monetary Fund put global private credit assets at about
\$2.1 trillion in 2023, counting assets and capital committed but not yet invested, mostly in North
America and Europe \parencite{imf2024gfsr}. This section takes the contract apart, prices its pieces,
models its covenant, and then reads the accounts of listed lenders to see what the reported returns\ix{Reported returns} show
and what they smooth away. Results are tagged \kind{exact} for a closed form, \kind{simulation} for a Monte
Carlo\ix{Monte Carlo simulation} estimate and \kind{historical} for a computation on recorded data.

\subsection{Who lends, and in what vehicle}

A private credit\ix{Private credit} fund is a closed-end\ix{Drawdown fund} pooled vehicle that originates or invests in loans\ix{Loan} to private
businesses and need not register as an investment company; like the funds of Chapter~2 it calls the
capital its investors\ix{Investor} have committed as it lends \parencite{cai2024}. A business
development company\ix{Business development company} (BDC) is a closed-end investment company created by the Small Business Investment
Incentive Act of 1980 to channel capital to small and mid-sized private businesses \parencite{cai2024}.
Some BDCs are listed and trade like shares\ix{Stock}; others are not listed and offer to buy back shares each
quarter. A BDC must keep asset coverage\ix{Asset coverage ratio} of senior securities of at least 150\% once its shareholders or
board have approved that level, against 200\% before, so it can borrow up to two dollars for every
dollar of net assets \parencite{arcc2025}. Ares Capital reported asset coverage of
\dlCoverArcc\% at the end of 2025, or debt of about \dlDeArcc{} times net assets, and lends to
middle-market companies, which it defines as those with annual EBITDA between \$10 million and \$250
million \parencite{arcc2025}.

The Bank of Canada finds that private credit is a limited and stable source of funds for Canadian
companies: loans from non-banks have stayed at about 15\% of the loans to Canadian businesses for a decade
\parencite{boc2026private}. Canadian institutions are large lenders abroad, mostly in the United States. The
Bank estimates the private lending of Canadian investors together with bank loans to private credit funds
at about \$500 billion early in 2026. The three largest life insurers\ix{Insurance} held just over \$200 billion, about
22\% of their invested assets, the large pension funds\ix{Pension plan} an estimated \$215 billion at the end of 2025, about
9\%, investment funds \$54 billion, and the banks had lent at least \$40 billion to managers of private
credit funds, about 1\% of their lending. Its Financial Stability Report notes that the sector is opaque,
increasingly linked to banks, and untested in a downturn \parencite{bocfsr2026}.

\subsection{The unitranche and its two lenders}

A \emph{unitranche\ix{Unitranche}} loan combines senior and subordinated debt in one loan, generally with a first lien\ix{Senior secured loan}\ix{Lien},
at a single blended rate \parencite{arcc2025}. The borrower sees one lender and one document. Behind it the
lenders may agree among themselves, in an intercreditor agreement, that one of them is repaid first: the
\emph{first out} piece is paid before the \emph{last out} piece, which earns more for standing behind it.
The borrower's rate does not change; the split only reallocates the loss and the spread.

\begin{proposition}[Splitting a unitranche]\label{prop6:unitranche}
Let a unitranche of face $F$ be split into a first out piece of face $wF$ and a last out piece of face
$(1-w)F$, paid from the recovery $Y$ at the horizon $T$. Then
\begin{enumerate}[label=(\roman*)]
  \item the expected losses\ix{Expected loss} add: $\mathrm{EL}_U=\mathrm{EL}_F+\mathrm{EL}_L$, with
        $\mathrm{EL}_F=\Pi(wF)$, $\mathrm{EL}_L=\Pi(F)-\Pi(wF)$ and $\Pi$ the put of \cref{prop6:slice};
  \item if the pieces pay spreads $s_F$ and $s_L$ on their faces, the borrower's spread $s_U$ is the
        blend $s_U=w\,s_F+(1-w)\,s_L$, so $s_L=(s_U-w\,s_F)/(1-w)$;
  \item with loss spreads $\ell_i=-\ln(1-\mathrm{EL}_i/\text{face}_i)/T$, which repay the expected loss of
        each piece over $T$ years, the split that gives both pieces the same premium $\pi$ over their
        loss spreads is $s_i=\pi+\ell_i$ with $\pi=s_U-w\,\ell_F-(1-w)\,\ell_L$.
\end{enumerate}
\end{proposition}

\begin{proof}
(i) is \cref{prop6:slice} with $b_{\mathrm{lo}}=0,b_{\mathrm{hi}}=wF$ and with $b_{\mathrm{lo}}=wF,b_{\mathrm{hi}}=F$,
and the sum telescopes to $\Pi(F)$. (ii) The interest paid on the pieces, $s_FwF+s_L(1-w)F$, must equal
the interest $s_UF$ paid by the borrower. (iii) Substituting $s_i=\pi+\ell_i$ into (ii) and solving for
$\pi$.
\end{proof}

The blend in (ii) is exact in the spreads. The loss spreads are not linear in the expected losses, so
the loss spread of the whole loan is not the blend of those of its pieces, and (iii) rather than the loss
spread of the whole loan is what fixes the premium. Take an illustration. A company has an enterprise
value\ix{Enterprise value} of \utV, lognormal at $T=\utT$ years with no expected growth\ix{Expected growth} and a log volatility of $\utSig\%$ a
year, and a sale after default costs \utCost\% of it. It borrows a unitranche of \utFace{} split \utFo{}
first out and \utLo{} last out, so $w=\utW\%$. \kind{exact} The five year expected losses are \utElU\% of
face for the whole loan, \utElF\% for the first out piece and \utElL\% for the last out piece, and the loss
spreads are \utLsU, \utLsF{} and \utLsL{} basis points\ix{Basis point}; the blend of the pieces' loss spreads is
\utLsBlend{} basis points, more than the whole loan's. Set the borrower's spread to \utSu{} basis points, the median
income yield less the bill rate of the listed BDCs\ix{Business development company} below in \dlLast. Then (iii) gives a premium of \utPi{} basis points to each
lender, a first out spread of \utSf{} and a last out spread of \utSl. The first out lender gives up most
of the spread for a loss that is a small fraction of the last out lender's. The enterprise value and its
volatility are the inputs that matter, and in practice both are estimates.

\subsection{Spreads, default and recovery}

\Cref{sec6:hazard} gave the spread of a credit default\ix{Default} swap with a constant hazard rate\ix{Hazard rate} $\lambda$ and
recovery\ix{Recovery rate} $R$ as about $\lambda(1-R)$. The same relation, read the other way, gives the default rate\ix{Default rate} a
lender's loss spread can absorb: a loss spread of $s$ pays for an annual default rate of $s/(1-R)$.
\textcite{cai2024} report, from KBRA DLD, recoveries after default of about \dlRecDirect\% for direct
loans\ix{Loan} against \dlRecSynd\% for broadly syndicated loans\ix{Leveraged loan} and \dlRecHy\% for high-yield bonds, and they
attribute the low recovery of direct loans to lending to sectors with few tangible assets, such as
software and health care services. \kind{exact} At those recoveries 100 basis points of loss spread
covers an annual default rate of \dlPdPerDirect\% on direct loans, \dlPdPerSynd\% on syndicated loans
and \dlPdPerHy\% on high-yield bonds, and a default rate of 2\% a year costs \dlLossDirectTwo\% a year on
direct loans against \dlLossSyndTwo\% on syndicated loans. The same paper reports a mean interest
coverage, EBITDA over interest, of about \dlIcrPrivate{} for private credit borrowers against
\dlIcrLev{} for leveraged loan borrowers. A direct lender\ix{Direct lending} is paid a higher spread for a borrower that is
more levered, recovers less when it fails, and has no market price to say how it is doing.

Part of a direct lender's return comes before the first coupon\ix{Coupon}. A loan is often funded below par, at
an original issue discount\ix{Original issue discount}, or with an upfront fee, and \textcite{cai2024} note that private loans can
carry high prepayment\ix{Prepayment} penalties. Both raise the yield, and by more when the loan is repaid early.

\begin{proposition}[Yield to life]\label{prop6:oid}
A loan bought at $1-u$ pays a coupon $c$ continuously and par at $T$. Its continuously compounded yield
$y$ solves $1-u=c\,(1-e^{-yT})/y+e^{-yT}$, and
\[
  y=c+\frac{c}{1-e^{-cT}}\,u+O(u^2),\qquad \frac{c}{1-e^{-cT}}>\frac1T ,
\]
so the discount adds a little more than $u/T$ a year, and a shorter life raises the yield.
\end{proposition}

\begin{proof}
The price is the integral of $ce^{-yt}$ over $[0,T]$ plus $e^{-yT}$. At $y=c$ it is exactly one, and
its derivative in $y$ there is $-\int_0^Tt\,ce^{-ct}\,dt-Te^{-cT}=-(1-e^{-cT})/c$, which is
smaller than $T$ in absolute value because $1-e^{-x}<x$ for $x>0$. Solving the first-order expansion
for $y-c$ gives the result.
\end{proof}

\kind{exact} A loan paying \oiC\% bought at a discount of \oiU{} points yields \oiBpFive{} basis points more
than its coupon if it lasts five years, against \oiApproxFive{} from $u/T$, and \oiBpThree{} more if it is
repaid in three, against \oiApproxThree; the first-order term of the proposition gives \oiFirstFive{}
and \oiFirstThree. A shorter life spreads the discount over fewer years.

\subsection{Floating rates and interest coverage}

A floating-rate\ix{Floating rate} loan moves the risk of the base rate from the lender to the borrower. If EBITDA\ix{EBITDA} is $X$
and debt $D$ pays the base rate $r$ plus a spread $s$, interest coverage\ix{Interest coverage} is $X/\bigl(D(r+s)\bigr)$, and
a rise in the base rate from $r_0$ to $r_1$ divides it by $(r_1+s)/(r_0+s)$ with EBITDA unchanged.
\kind{historical} The one month US bill averaged \icBillLow\% in 2021 and reached \icBillHigh\% in
\dlBillHighQ. With the spread of \utSu{} basis points used above, the cost of the same debt rose by
\icCostRise\%. A borrower with the mean coverage of about \dlIcrPrivate{} reported by \textcite{cai2024}
at the higher rates would, with the same EBITDA, have had coverage of about \icLow{} at the rates of
2021. Coverage of one is the point at which EBITDA just pays interest, the trigger of the covenant-lite
loan below. The lenders' income rose with the rate, as \cref{fig6:direct} shows; the borrowers' margin
of safety fell by the same mechanism.

\subsection{The covenant as a barrier}

A maintenance covenant\ix{Covenant} is tested every quarter whether or not the borrower has missed a payment: debt
over trailing EBITDA\ix{EBITDA} must stay below a level written into the loan. A breach lets the lender reprice,
take more security or force a sale while the business still has value. A loan without such tests, a
covenant-lite loan, leaves the lender waiting until the borrower cannot pay. Industry commentary cited by
\textcite{cai2024} suggests that recent large private loans carry fewer maintenance covenants as direct
lenders\ix{Direct lending} compete with banks.

\begin{proposition}[Covenant breach]\label{prop6:covenant}
Let EBITDA follow $X_t=X_0e^{\nu t+\sigma W_t}$ and let debt $D$ be fixed. A covenant $D/X_t\le c$, tested
continuously, is breached by $T$ with probability
\[
  P\Bigl(\min_{t\le T}X_t\le b\Bigr)=\Phi\Bigl(\frac{a-\nu T}{\sigma\sqrt T}\Bigr)
  +e^{2\nu a/\sigma^2}\,\Phi\Bigl(\frac{a+\nu T}{\sigma\sqrt T}\Bigr),\qquad b=\frac Dc,\quad a=\ln\frac b{X_0}<0 .
\]
If the business is sold at the first breach for a multiple $m$ of EBITDA, after costs $\theta$, the
lender recovers $\min\bigl(1,(1-\theta)m/c\bigr)$ of its claim.
\end{proposition}

\begin{proof}
The first part is the law of the minimum of a Brownian motion with drift, which follows from the
reflection principle used in Chapter~9 for fund closures \parencite{shreve2004}. At the first breach
$X=b=D/c$, so a sale raises $(1-\theta)mD/c$ against a claim of $D$.
\end{proof}

Take a borrower with debt of \cvLev{} times EBITDA, no expected growth\ix{Expected growth} and an EBITDA volatility of
\cvSig\% a year. \kind{exact} A covenant at \cvCov{} times trips when EBITDA has fallen \cvDrop\%, with
probability \cvPcov\% within \cvT{} years. A covenant-lite loan at a rate of \cvRate\% fails only when
EBITDA no longer covers interest, a fall of \cvDropLite\%, with probability \cvPlite\%. Selling at \cvMult{}
times EBITDA after costs of \cvCost\%, the lender recovers \cvRecCov\% at the covenant and \cvRecLite\%
at the payment default. The covenant trips \cvTimes{} times as often, and at a level where a sale would
still repay almost all of the loan, which is what gives the lender room to negotiate. \Cref{fig6:covenant} shows how the breach probability falls as the covenant
is loosened. \kind{simulation} With quarterly tests in mind the formula can be checked by simulating
monthly paths: 200{,}000 paths monitored monthly breach in \cvMcRaw\% of cases, fewer than the continuous
\cvPcovOne\% because a path can dip below the level between tests, and with the barrier shifted by the
correction of \textcite{broadie1997} the count is \cvMc\%, with a standard error of \cvMcSe{} points.

\begin{figure}[t]
  \centering
  \includegraphics[width=0.96\textwidth]{ch06/fig14_covenant}
  \caption{\kind{exact} Probability that a maintenance covenant on debt over EBITDA is breached within
  \cvT{} years by \cref{prop6:covenant}, for a borrower with opening leverage of 5, no expected growth in
  EBITDA and three volatilities. An illustration.}
  \label{fig6:covenant}
\end{figure}

\subsection{What a listed lender charges}

A BDC is managed by an adviser that is paid a base fee on assets and an incentive fee on income. Ares
Capital pays a base fee of \dlBaseLo\% a year on total assets other than cash financed up to a debt to
equity ratio of one, and \dlBaseHi\% on assets financed above it. Its income based fee is nothing while
pre-fee net investment income is below a hurdle\ix{Hurdle rate} of \dlHurdleQ\% of net assets a quarter, all of the
income between the hurdle and \dlCatchQ\%, the catch-up\ix{Catch-up}, and \dlIncRate\% of all of it above that; a
further fee is charged on realised capital gains \parencite{arcc2025}. The catch-up ends where
$\dlIncRate\%$ of the whole equals the excess over the hurdle, $1.75/(1-0.2)=2.1875$, as for the soft
hurdle with full catch-up of Chapter~9.

\begin{proposition}[A base fee charged on gross assets]\label{prop6:bdcfee}
A fund with debt to equity $d$ and no cash, paying $f_1$ a year on assets financed up to $d=1$ and $f_2$
above, pays base fees of $f_1(1+\min(d,1))+f_2(d-1)^+$ of net assets a year. If its assets yield $y$ and
its debt costs $r_D$, its return on net assets before fees is $y(1+d)-r_Dd$.
\end{proposition}

\begin{proof}
Each unit of net assets carries $1+d$ of assets, of which up to $1+\min(d,1)$ are charged $f_1$ and
$(d-1)^+$ are charged $f_2$; the assets earn $y(1+d)$ and the debt costs $r_Dd$.
\end{proof}

\kind{exact} At Ares Capital's leverage\ix{Leverage} of \dlDeArcc{} the base fee is \dlBaseNaArcc\% of net assets a year,
and at the legal limit of $d=\dlDeMax$ it would be \dlBaseNaMax\%. With its reported yield of
\dlYieldArcc\% on debt investments at amortised cost at the end of 2025 \parencite{arcc2025} and, as an
assumption, debt at \dlExRd\%, the return on net assets before fees is \dlExGross\%, the base fee takes
\dlExBase\% and the income fee \dlExInc\%, leaving \dlExNet\% before other expenses and credit losses\ix{Loss rate}:
fees take \dlExFeeShare\% of the return that leverage produces. The base fee on gross assets pays the
adviser for leverage whether or not leverage pays the shareholders.

Leverage also sets how much loss the asset coverage\ix{Asset coverage ratio} test can absorb.

\begin{proposition}[Losses that break the coverage test]\label{prop6:coverage}
A fund with debt to equity $d$ and no other liabilities has asset coverage $(1+d)/d$. It keeps coverage
of at least $\kappa$ after losing a share $L$ of its assets if and only if $L\le1-\kappa d/(1+d)$, and the
loss costs its shareholders $(1+d)L$ of net assets. If the loans\ix{Loan} default\ix{Default} as in the large pool\ix{Large homogeneous pool} of
\cref{prop6:lhp}, with probability\ix{Probability} $p$, correlation $\rho$ and loss given default\ix{Loss given default} $\mathrm{LGD}$, the
probability that defaults over the period break the test is
\[
  1-\Phi\Bigl(\frac{\sqrt{1-\rho}\,\Phi^{-1}(x^*)-\Phi^{-1}(p)}{\sqrt\rho}\Bigr),\qquad
  x^*=\frac{1-\kappa d/(1+d)}{\mathrm{LGD}},
\]
when $0<x^*<1$.
\end{proposition}

\begin{proof}
After the loss, assets are $(1+d)(1-L)$ against debt $d$, so coverage is at least $\kappa$ exactly when
$L\le1-\kappa d/(1+d)$, and the loss of $(1+d)L$ falls on equity of one. The loss share is
$\mathrm{LGD}$ times the default share, so the test breaks when the default share exceeds $x^*$, whose
probability is one less \cref{prop6:lhp} at $x^*$.
\end{proof}

\kind{exact} At Ares Capital's leverage the test allows a loss of \acThrArcc\% of assets, at $d=1$ a loss of
\acThrOne\%, and at $d=1.5$ only \acThrOneHalf\%. Take as an illustration a one year default
probability of \acP\% and a correlation of \acRho, with the loss given default of direct loans from
\textcite{cai2024}. Defaults alone then break the test within a year with probability \acBreachArcc\% at
Ares Capital's leverage and \acBreachOneHalf\% at $d=1.5$. In the worst year in a thousand of this model
\acQ\% of the loans default, the book loses \acQloss\% of its assets, and net assets fall by \acQeq\% at
Ares Capital's leverage and by \acQeqTwo\% at the legal limit of two. Losses on marks, which come before
defaults, would break the test sooner.

\subsection{What the accounts of nineteen lenders show}

Listed BDCs\ix{Business development company} file quarterly accounts with the Securities and Exchange Commission, tagged in XBRL, and the
SEC publishes the tagged figures \parencite{secxbrl}. \texttt{scripts/pull\_bdc.py} reads them for
\dlFirms{} listed BDCs, from Ares Capital to Gladstone Investment, and builds a table of their net
assets, investments at fair value\ix{Fair value} and at amortised cost, investment income, fees, liabilities and the
net increase in net assets resulting from operations, which is net investment income plus realised and
unrealised gains. The tagging is broad from 2021, and for investment income from late 2022, so the
results below run from \dlFirst{} to \dlLast{}. \Cref{tab6:bdc} lists the companies and
\cref{fig6:direct} draws the averages. All are \kind{historical}.

\begin{table}[t]
  \centering
  \small
  \caption{\kind{historical} Nineteen listed business development companies in \dlLast, from their XBRL
  filings with the SEC: net assets, total liabilities over net assets, investments at fair value as a
  percentage of amortised cost, the median share of total investment income paid in management and
  incentive fees (a dash where the fees are not tagged), and the average annualised quarterly return on
  net assets from \dlFirst. Ordered by net assets.}
  \label{tab6:bdc}
  % Generated by code/ch06/direct_lending.py. Do not edit by hand.
\begin{tabular}{@{}lrrrrr@{}}
\toprule
company & net assets (\$bn) & liab./NA & FV/cost (\%) & fees (\%) & return (\%)\\
\midrule
Ares Capital & 13.9 & 1.20 & 98.9 & 26 & 10.8\\
Blue Owl Capital & 7.0 & 1.18 & 99.0 & 23 & 9.3\\
Blackstone Secured Lending & 5.9 & 1.32 & 96.4 & 19 & 10.4\\
FS KKR Capital & 5.1 & 1.31 & 90.6 & 23 & 3.6\\
Golub Capital BDC & 3.7 & 1.25 & 97.8 & 25 & 7.8\\
Main Street Capital & 3.2 & 0.88 & 115.5 & -- & 17.6\\
Hercules Capital & 2.3 & 1.06 & 99.6 & -- & 14.6\\
Sixth Street Specialty Lending & 1.5 & 1.29 & 99.6 & 21 & 11.5\\
Oaktree Specialty Lending & 1.4 & 1.07 & 91.5 & 18 & 3.8\\
Goldman Sachs BDC & 1.4 & 1.42 & 93.7 & -- & 6.6\\
Barings BDC & 1.1 & 1.26 & 98.3 & 22 & 7.4\\
Carlyle Secured Lending & 1.1 & 1.26 & 97.7 & 22 & 8.7\\
Capital Southwest & 1.1 & 1.19 & 98.6 & -- & 10.4\\
New Mountain Finance & 1.0 & 1.34 & 92.9 & 22 & 5.2\\
PennantPark Floating Rate & 1.0 & 1.59 & 95.3 & 19 & 6.4\\
SLR Investment & 1.0 & 1.61 & 98.1 & -- & 7.2\\
Fidus Investment & 0.7 & 1.03 & 102.0 & 25 & 8.8\\
Gladstone Investment & 0.6 & 1.01 & 122.7 & 27 & 17.2\\
BlackRock TCP Capital & 0.6 & 1.67 & 89.9 & -- & $-$2.4\\
\bottomrule
\end{tabular}

\end{table}

\emph{Floating rates\ix{Floating rate}.} The median\ix{Median} ratio of total investment income to average amortised cost rose from
\dlYieldFirst\% in \dlYieldFirstQ{} to \dlYieldPeak\% in \dlYieldPeakQ{} and was \dlYieldLast\% in \dlLast,
while the one month US bill averaged \dlBillFirst\%, then \dlBillPeak\% at its peak, then \dlBillLast\%. A
regression of the median yield on the bill rate over the \dlYieldN{} quarters has a slope of \dlPass{},
with a standard error of \dlPassSe{} and an $R^2$ of \dlPassRsq: a change in the base rate passes into
the lenders' income about one for one, as floating-rate loans require. The median income yield less the
bill rate was \dlSpreadLastBp{} basis points in \dlLast, which includes fees, discounts earned and
dividends as well as the contractual spread.

\emph{Leverage and fees.} In \dlLast{} the median ratio of total liabilities to net assets was
\dlLevMed, and the highest \dlLevMax, within the limit of two. The management and incentive fees took a
median \dlFeeShare\% of investment income, or \dlFeeNa\% of net assets a year. The companies differ more
than the medians suggest: two carry their investments well above cost and three about ten points below
it (\cref{tab6:bdc}).

\emph{Marks.} Investments at fair value were \dlMarkFirst\% of their cost across the BDCs in \dlFirst{}
and \dlMarkLast\% in \dlLast, the lowest of the period. A fall in the ratio is the lenders' own estimate
of what their loans have lost, in default or in value, since they were made. Ares Capital reported loans
on non-accrual\ix{Nonaccrual} status, those no longer paying, at \dlNonaccCost\% of its investments at amortised cost
and \dlNonaccFv\% at fair value at the end of 2025 \parencite{arcc2025}.

\begin{figure}[t]
  \centering
  \includegraphics[width=0.96\textwidth]{ch06/fig13_direct_lending}
  \caption{\kind{historical} Nineteen listed business development companies, from their quarterly filings
  with the SEC. (a) The median ratio of total investment income to average investments at amortised
  cost, annualised, and the US one month bill rate. (b) Investments at fair value less cost, summed across
  the companies, as a share of cost, and the equally weighted average quarterly return on net assets,
  the result of operations over the previous quarter's net assets.}
  \label{fig6:direct}
\end{figure}

\subsection{Stale marks and de-smoothing}

A loan with no market price is marked by its lender, usually each quarter, with a model and comparable
deals. The IMF lists stale\ix{Stale pricing} and potentially subjective valuations\ix{Valuation} among the vulnerabilities of private
credit, together with semi-liquid vehicles and layers of leverage, and warns that a severe downturn could bring a
delayed realisation of losses followed by a spike in defaults and large valuation markdowns
\parencite{imf2024gfsr}. Chapter~8 gave the tools.
If reported returns\ix{Reported returns} follow $r^*_t=(1-\phi)r_t+\phi r^*_{t-1}$, the smoothing\ix{Return smoothing} proposition of Chapter~8 says the reported
volatility\ix{Measured volatility} understates the true one by $\sqrt{(1-\phi)/(1+\phi)}$, the first autocorrelation\ix{Autocorrelation}\ix{First-order autocorrelation} of reported
returns estimates $\phi$, and $r_t=(r^*_t-\phi r^*_{t-1})/(1-\phi)$ recovers the true return. And if marks
lag the market, a regression on the market's return this quarter and last \parencite{dimson1979} gives a
lagged beta\ix{Dimson beta}.

\kind{historical} The average quarterly return on net assets of the BDCs\ix{Business development company} was \dlRoeMean\% a year with a
volatility of only \dlRoeVol\%, and its worst quarter, \dlRoeMinQ, was \dlRoeMin\%. Its first
autocorrelation over the \dlRoeN{} quarters is \dlRhoEw. Pooled within companies, over \dlRhoN{} pairs of
quarters, the first autocorrelation is \dlRho{} with a standard error of \dlRhoSe, so the persistence is
in the average more than in each lender, which is what marks that follow common market inputs with a
delay would produce. On the US equity market,
compounded to quarters from the French data, the average has a beta of \dlBetaZero{} on the current
quarter and \dlBetaOne{} on the last, with a $t$ statistic\ix{Test statistic} of \dlBetaOneT{} on the lag over \dlDimN{}
quarters: the lenders' marks follow the market with a delay.

The sample is short and has no recession\ix{Recession} in it. The Form PF statistics of Chapter~9 give the credit
hedge funds\ix{Hedge fund} of large advisers from 2013 \parencite{ofrhfm}. Their quarterly net return has a first
autocorrelation of $\phi=\dlPfPhi$ over \dlPfN{} quarters; the same proposition then scales the reported
volatility of \dlPfVol\% by \dlPfRatio, and de-smoothing\ix{Unsmoothing} the series gives a volatility of \dlPfVolDs\%.
The worst quarter, \dlPfMinQ, was \dlPfMin\% as reported and \dlPfMinDs\% de-smoothed. The correction is
modest for these funds. An autocorrelation measured in years without a recession may also understate the
smoothing that a long downturn would reveal, since the proposition assumes a constant $\phi$.

\subsection{Liquidity: the tender queue}

A non-listed BDC offers liquidity\ix{Liquidity risk} by tender. Blackstone Private Credit Fund, for example, may buy back
in each quarter up to \dlCap\% of its shares\ix{Stock} by net asset value\ix{Net asset value}\ix{Asset value}, at the quarter's net asset value per
share, or \dlEarly\% of it for shares held less than a year, and when requests exceed the offer only part
of each request is filled \parencite{bcred2025}. The fund offers investors\ix{Investor} at most \dlCap\% a quarter while
it holds loans that, at Ares Capital for example, generally run three to ten years \parencite{arcc2025}.

\begin{proposition}[The tender queue]\label{prop6:tender}
A fund buys back a share $c$ of its net assets each quarter, pro rata among tenders that total $f>c$. An
investor who tenders all its remaining holding every quarter still holds $(1-c/f)^n$ of it after $n$
quarters, and needs $\ln 0.1/\ln(1-c/f)$ quarters to get nine tenths out.
\end{proposition}

\begin{proof}
Each quarter a share $c/f$ of every tender is filled, so the remaining holding is multiplied by
$1-c/f$.
\end{proof}

\kind{exact} If a tenth of the fund asks to leave each quarter, half of each request is filled and nine
tenths of a holding take \dlQuartersTen{} quarters to leave. If a fifth asks, a quarter is filled, nine
tenths take \dlQuartersTwenty{} quarters, and after a year \dlLeftFourTwenty\% of the holding is still in.
The queue protects the remaining investors from forced sales\ix{Liquidation}, and it is also why the first investors to
ask fare best.

\FloatBarrier
% =====================================================================================
\section{Case study: could a five year commercial mortgage have been refinanced?}
\label{sec6:case}

\kind{historical and simulation} A typical Canadian commercial mortgage\ix{Mortgage}\ix{Commercial mortgage} has a five year term and a much longer amortisation\ix{Amortisation}, so most
of the principal is still owed at maturity\ix{Maturity} and has to be refinanced. Whether that is possible depends
on the rate at maturity, the income the property earns by then, and the value a new lender will put
on it. The study takes a stylised loan written at every month-end from \csFirst{} to \csLast, the
\csN{} vintages\ix{Vintage year} whose maturity falls within the data, and asks whether a new lender would have lent
enough to repay it.

\paragraph{The loan}
At origination the property is worth 100. Its cap rate\ix{Capitalisation rate} depends on the rates of the day: it is
\csCap\% when the five year Government of Canada\ix{Government of Canada bond} yield is \csYRef\% and moves by $\beta=\csBeta$ points
for each point of yield, $q_0=\csCap\%+\beta(y_0-\csYRef\%)$, so income is $\mathcal N_0=100\,q_0$.
Over the vintages $q_0$ runs from \csCapZeroMin\% to \csCapZeroMax\%. The mortgage rate is $y_0$ plus
\csSpread{} points, stated as a yearly rate compounded half-yearly\ix{Compounding!semi-annual}, as Canadian mortgage rates are,
and paid monthly over a \csAmort{} year amortisation, so the monthly rate is $(1+j/2)^{1/6}-1$ for a
quoted rate\ix{Quoted rate} $j$. The lender applies two tests at origination and lends the smaller of \csLtv\% of value
and the amount whose debt service coverage\ix{Debt service coverage ratio} is \csDscrMin, $\mathcal N_0/(\csDscrMin\,k_0)$, with $k_0$
the annual payment per unit of loan. Coverage sets the loan in \csDscrBindsZero{} of the \csN{}
vintages, and loan to value\ix{Loan to value} at origination runs from \csLtvZeroMin\% to \csLtvZeroMax\%. Without the
coverage test, a loan of \csLtv\% at a fixed \csCap\% cap rate would start below a coverage of
\csDscrMin{} whenever the yield is above \csOldYCut\%, which is the case for \csOldBelow{} vintages; no
lender would write those loans\ix{Loan}, and a refinancing\ix{Refinancing risk} test that the original loan could not pass says
little about refinancing risk.

The loan amortises, so its balance at maturity is $b$. Over the term income grows by \csG\% a year,
so $\mathcal N_1=\mathcal N_0(1.02)^5$. Cap rates keep moving with yields,
$q_1=q_0+\beta(y_1-y_0)$, where $y_1$ is the five year yield at maturity, so the property is worth
$V_1=\mathcal N_1/q_1$. A new lender quotes the yield then plus the same spread and applies the same
two tests: it lends the smaller of \csLtv\% of $V_1$ and $\mathcal N_1/(\csDscrMin\,k_1)$, with $k_1$ the
new annual payment per unit of loan. The \emph{cover} $\gamma$ is that amount divided by $b$, and the
loan refinances if $\gamma\ge1$. The data are the month-end five year yields from the file of
Chapter~4, which runs to \csMatLast. Reading the quoted rates as compounded monthly instead would raise
the effective rate by at most \csConvMax{} basis points\ix{Basis point} and change the number of vintages that fail the
test without shocks from \csFails{} to \csMonthlyFails.

\paragraph{Without shocks}
On these assumptions \csFails{} of the \csN{} vintages have a cover below 1, all written between
\csFailFirst{} and \csFailLast, when the five year yield was at most \csFailYMax\%. The median cover is
\csCoverMed, and the coverage test sets the new loan in \csDscrBindsOne{} vintages. When coverage
sets both the original and the new loan, the cap rate drops out of the cover: the loan is proportional
to $\mathcal N_0$, and so are $b$ and $\mathcal N_1$, so $\gamma$ depends only on the two loan constants
and on income growth\ix{Rent growth}. A loan written when rates were low and refinanced when they had risen is then
the one at risk. The lowest cover, \csCoverMin, belongs to the vintage of \csCoverMinDate. The last
vintage, \csLast, was written at a yield of \csLastYZero\% and a rate of \csRateZero\%, with a cap rate
of \csLastCapZero\%, so coverage limited the loan to \csLastLoan{} with a coverage of \csLastDscrZero.
By maturity the yield was \csLastYOne\% and the annual payment per 100 of loan had gone from
\csLastKZero{} to \csLastKOne. The balance was \csLastB, and a new loan equal to it would have had a
coverage of \csLastDscrOne, just short of \csDscrMin, while its loan to value at the new cap rate of
\csLastCapOne\% was \csLastLtvOne\%, inside the \csLtv\% limit. The cover is \csLastCover, and what stops the
refinancing is the higher rate. \Cref{fig6:refi} shows the yield and the cover for every vintage.

\begin{figure}[t]
  \centering
  \includegraphics[width=0.92\textwidth]{ch06/fig11_refinance}
  \caption{\kind{historical} Top: month-end five year Government of Canada yield, \csFirst{} to \csMatLast. Bottom: the
  cover $\gamma$ of each vintage by origination month (left axis, solid), with the line at 1, and the
  probability that the loan cannot refinance once income is uncertain (right axis, dashed).}
  \label{fig6:refi}
\end{figure}

\paragraph{With shocks}
Income at maturity is not known at origination. Let both income and value at maturity be multiplied
by $\zeta=\exp\bigl(\sigma_MM+\sigma_\varepsilon\varepsilon-\tfrac12\sigma^2\bigr)$ with
$\sigma^2=\sigma_M^2+\sigma_\varepsilon^2$, where $M$ is common to every property and $\varepsilon$ is
specific to one, both standard normal, so that $\E\zeta=1$.

\begin{proposition}[Refinancing as a one-factor model]\label{prop6:refi}
A loan of a vintage with cover $\gamma$ fails to refinance exactly when $\zeta<1/\gamma$, which has
probability\ix{Probability}
\[
  p=\Phi\Bigl(\frac{\ln(1/\gamma)+\tfrac12\sigma^2}{\sigma}\Bigr).
\]
For loans of the same vintage with independent $\varepsilon_i$ and a common $M$, failures follow
\cref{def6:copula} with this $p$ and $\rho=\sigma_M^2/\sigma^2$.
\end{proposition}

\begin{proof}
Both limits scale with $\zeta$: $\csLtv\%\cdot V_1\,\zeta$ because $V_1\,\zeta=\mathcal N_1\,\zeta/q_1$, and
$\mathcal N_1\,\zeta/(\csDscrMin\,k_1)$ directly, since $q_1$ and $k_1$ do not depend on $\zeta$. So the largest new
loan is $\gamma b\zeta$, and the loan fails when $\gamma \zeta<1$. Writing
$Z=(\sigma_MM+\sigma_\varepsilon\varepsilon)/\sigma$, this is
$Z<\bigl(\ln(1/\gamma)+\tfrac12\sigma^2\bigr)/\sigma=\Phi^{-1}(p)$, and
$Z_i=\sqrt\rho\,M+\sqrt{1-\rho}\,\varepsilon_i$ with $\sqrt\rho=\sigma_M/\sigma$, which is the latent variable $X_i$ of \cref{def6:copula}.
\end{proof}

With $\sigma_M=\csSigM\%$ and $\sigma_\varepsilon=\csSigE\%$, $\sigma=\csSig\%$ and $\rho=\csRho$. The
dashed line in \cref{fig6:refi} is $p$ for each vintage. Its median over all vintages is
\csPdMed\%, and over the vintages written before 2020 it is \csPdPreMed\%. The largest, \csPdMax\%,
belongs to the vintage of \csPdMaxDate, which fails even without a shock. The answer depends on how
cap rates follow yields, but only through which test binds. For the last vintage, $\beta=1$ leaves the
cover at \csCoverBetaOne{} and the failure probability at \csPdBetaOne\%, since coverage still sets
both loans. With $\beta=0$ every cap rate is \csCap\%, value sets the original loan at \csLtv\%, and
the cover is \csCoverBetaZero{} with a failure probability of \csPdBetaZero\%.

\paragraph{From failure to loss}
A loan that cannot refinance is assumed to be sold, the lender receiving $\eta\zeta$ with
$\eta=(1-\csCost\%)V_1$ against the balance $b$. The lender's loss is $(b-\eta\zeta)^+$; sale
proceeds exceeding the debt do not produce a gain on the loan. With $u=\min(1/\gamma,\,b/\eta)$ and
$z=(\ln u+\tfrac12\sigma^2)/\sigma$, the expected loss\ix{Expected loss} is
\[
  \E\bigl[(b-\eta\zeta)^+\mathbf 1\{\zeta<1/\gamma\}\bigr]=b\,\Phi(z)-\eta\,\Phi(z-\sigma),
\]
because the loss is $b-\eta\zeta$ exactly when $\zeta<u$, $\mathbb P(\zeta<u)=\Phi(z)$ and $\E[\zeta\mathbf 1\{\zeta<u\}]=\Phi(z-\sigma)$.
For the last vintage this is \csEl{} basis points of the balance. The probability of failure is high
and the loss given failure small, because the shortfall that stops a refinancing is a few points of
value while a sale at a loan to value near \csLastLtvOne\% leaves a cushion. A simulation of
\csMcScen{} scenarios\ix{Scenario analysis} for a pool of \csLoans{} such loans (stream 7) gives an average loss of
\csMcEl{} basis points with standard error \csMcElSE, \csMcElZ{} standard errors from the formula.

The pool's number of failures has mean \csMcMeanFail{} in the simulation, with standard error
\csMcMeanFailSE, against \csMeanFail{} from \cref{prop6:refi}. Its distribution is wide:
\cref{prop6:cond} gives a probability of \csPoolTail\% that 25 or more of the \csLoans{} loans fail,
and the simulation \csPoolTailMc\% (\cref{fig6:pool}). A hypothetical claim that bears all pool losses above an attachment point of \csAttach\% of the
pool balance, up to the whole pool, so that its size is 99\% of the balance, has an expected loss of
\csSenEl{} basis points of the pool balance (\csSenElOwn{} of its own size), with standard error
\csSenElSE. It is hit in only \csSenPHit\% of scenarios, so it has its own simulation of
\csSenScen{} scenarios (stream 9). Each draws the common shock and \csLoans{} loan shocks as above;
a loan fails when $\zeta<1/\gamma$ with $\gamma=\csLastCover$, and then loses
$(b-\eta\zeta)^+/b$ with $b=\csLastB$ and $\eta=\csEta$, which is $(1-\csCost\%)$ of $V_1=\csLastVOne$.
If $\bar L$ is the average loss over the pool, the claim loses $\min\bigl((\bar L-1\%)^+,\,99\%\bigr)$
of the balance. The screen of \cref{sec6:private},
given this loan at origination with a sponsor rating of \csSponsor{} and no cap rate shift, scores it \csScore, a \csGrade; given the cap rate
shift that actually happened it scores \csScoreReal, still a \csGradeReal. Even given that shift, the
screen sees an exit loan to value of \csScreenExitLtv\%, far from the 90\% at which its exit score
reaches zero. It has no input for the new lender's tests, a loan to value of at most \csLtv\% and a
coverage of at least \csDscrMin, or for uncertainty\ix{Uncertainty} in income, and it is those, applied at the rate at
maturity, that moved the failure probability.

\begin{figure}[t]
  \centering
  \includegraphics[width=0.84\textwidth]{ch06/fig12_pool}
  \caption{\kind{exact and simulation} Number of loans out of \csLoans{} from the \csLast{} vintage that cannot refinance: the
  histogram of \csMcScen{} simulated scenarios and the distribution of \cref{prop6:cond} with
  $\rho=\csRho$, for every count from 0 to \csLoans.}
  \label{fig6:pool}
\end{figure}

% =====================================================================================
\section{What the models leave out}
\label{sec6:limits}

Every result in this chapter holds under hypotheses that real markets break, and the breaks matter
in different places.

The rate models are single-factor models fitted to one curve. In both models, every continuously compounded
zero-coupon yield is an affine function of the short rate\ix{Short rate} with a positive loading, so these yields move
in the same direction when the short rate changes; par yields also move in the same direction, although
they are not generally affine functions of the short rate. The models therefore cannot describe a twist of the curve; with one random
driver, level, slope and curvature cannot move independently, and the fitted Vasicek and CIR
parameters change from one date to the next, as \cref{tab6:fits} showed. Swaps\ix{Swap rate} are discounted here on
the government curve, while market swaps are discounted on the overnight curve and quoted at a spread
to government bonds\ix{Sovereign bond}. Caplet prices from a Gaussian model have a fixed skew, and a market smile
generally needs more parameters.

The Merton model\ix{Merton model} has one debt maturity\ix{Maturity}, continuous asset paths and no payouts, and its asset value\ix{Asset value}
is never observed. Firms have many debt issues, pay dividends and coupons, and can default before
their debt matures. The first-passage models that followed allow for that, but with a fixed barrier below today's asset value and continuous paths they too give short-dated
debt almost no spread; jumps and uncertain barriers are among the ways later models obtain one. The reduced-form model fixes the
recovery rate\ix{Recovery rate} and assumes that default is independent of interest rates\ix{Interest rate}, and a bootstrapped curve\ix{Bootstrapping a curve} is
only as good as the quotes it is built from: a step function can fit any set of quotes that admits
nonnegative hazards, which says nothing about whether the steps are where the risk is.

The Gaussian copula\ix{Copula} has one correlation\ix{Correlation} for every pair of names and thin joint tails. When tranche\ix{Tranche}
prices are inverted for a correlation, each tranche gives a different number, which is the model
telling its users that its shape is wrong. Dealers patched this by quoting a base correlation for each detachment point
of the traded indices, pricing every tranche as the difference of two equity tranches, while rating and risk models for bespoke pools often kept a single
correlation; neither repair changed the thin tails that mattered in 2007 and 2008. Homogeneous pools with equal probabilities and recoveries
simplify any real portfolio.

The screener's weights and thresholds are choices, and the grade reflects those choices as much as
the loan. In the case study the loan is stylised, the spread is the same in 2001 as in 2021, cap rates\ix{Capitalisation rate}
follow yields by a fixed rule, income grows at a fixed rate, nobody prepays, a lender who cannot
refinance sells at once with a fixed cost, and losses are counted at maturity without discounting or
lost interest. The shock to income is lognormal with one common factor. Each of these could be
relaxed in the code, and the dependence of the answer on $\beta$ is a warning that some of them
matter a great deal.

% =====================================================================================
\section{Summary}
\label{sec6:summary}

A bond's duration and convexity give its price change for a small parallel move in yield, with an
error of known sign, and matching the present value and duration of assets to a liability protects a
lender against such moves when the assets are more spread out in time than the liability. The Vasicek
and CIR models give zero-coupon prices in closed form and fit the Canadian curve to 5 to 10 basis
points in root mean square on the two dates tried. In the Ho-Lee model of \cref{prop6:holee} a futures
rate sits above the forward rate by $\tfrac12\sigma^2T_1T_2$, an amount that grows with the product of
the two dates. A swap rate on a single curve is a par yield. A caplet is a put on a zero-coupon bond, and
the fitted Gaussian short-rate model produces, in the examples shown, Black implied volatilities that fall with expiry and strike. In the Merton model
debt is a riskless bond less a put on the assets, short spreads vanish when the assets cover the
debt, and equity prices and volatility can be inverted for the assets. A hazard rate turns CDS quotes
into survival probabilities one interval at a time. A common factor makes defaults arrive together
without changing their mean; tranching redistributes the pool's expected loss and moves it from
equity to senior notes as losses spread out, and a structure whose notes outlast its collateral is much less exposed to forced sales, though not
immune to them. A lender's screening ratios are tied together by two identities and do not, by themselves,
identify the volatility of the collateral, which drives the expected loss of a slice. A unitranche
split into first out and last out pieces keeps the borrower's spread as the blend of the pieces'
spreads, and with the recoveries reported for direct loans a basis point of spread covers less default
than on syndicated loans. A maintenance covenant is a barrier on EBITDA that trips earlier and at a
better recovery than a payment default. The accounts of \dlFirms{} listed lenders show income that moved
with the base rate about one for one, fees of about a fifth of income, marks below cost in \dlLast{} and
returns on net assets that follow the equity market with a lag. Over \csN{} vintages
of a stylised Canadian commercial mortgage, sized at origination by both tests, \csFails{} vintages written
between \csFailFirst{} and \csFailLast{} fail the refinancing test even without shocks, and with uncertain
income the failure probability is largest, \csPdMax\%, for the vintage of \csPdMaxDate.

% =====================================================================================
\section{Exercises}
\label{sec6:exercises}

\begin{enumerate}[label=\textbf{6.\arabic*}, leftmargin=*, ref=6.\arabic*]
  \item\label{ex6:one} A five year bond pays a 4\% coupon twice a year and yields 5\%, compounded
  twice a year. Find its price, modified duration and convexity, and estimate its price after the
  yield rises by half a point with duration alone and with convexity. Compare with the exact price.
  \item\label{ex6:two} Check \cref{prop6:par} for a two year par bond at 4\% by computing the duration
  from the sum in \cref{def6:duration}.
  \item\label{ex6:three} A liability of 50 is due in six years and the rate is a flat 4\%, compounded
  continuously. How much must two zero-coupon bonds pay at one year and at ten years to match the
  liability's present value and duration? Can a parallel shift of the rate make the surplus negative?
  \item\label{ex6:four} In the Ho-Lee model with $\sigma$ of one percentage point, what is the
  convexity adjustment for a three month rate starting in three years?
  \item\label{ex6:five} A firm's assets are worth 120, with volatility 20\%, and its only debt is a
  zero-coupon bond with face 100 due in three years. The rate is 3\%. Find the values of equity and
  debt, the credit spread and the risk-neutral probability of default.
  \item\label{ex6:six} A five year CDS trades at 150 basis points and the recovery rate is 40\%. Use
  the credit triangle to estimate a flat hazard rate and the probability of surviving five years.
  \item\label{ex6:seven} In the large-pool limit with $p=2\%$ and $\rho=0.2$, find the default share that
  is exceeded with probability 1\%.
  \item\label{ex6:eight} A loan of 7 million against a property worth 10 million with net operating
  income of 600,000 is interest only at 7\%. Find its loan to value, debt yield and coverage, and check
  both identities of \cref{prop6:identities}.
  \item\label{ex6:nine} A loan of 1,000,000 at 4\%, compounded monthly, amortises over 25 years and is
  due in five. At maturity the rate is 6\%, the property earns 90,000 a year and a new lender requires a
  coverage of 1.25 on a 25 year amortisation. Find the balance due, the largest new loan and the cover.
  \item\label{ex6:ten} A unitranche pays a spread of 600 basis points and is split 40\% first out and
  60\% last out, with the first out lender paid 350 basis points. What does the last out lender earn?
  With the recovery of direct loans used in \cref{sec6:direct}, what annual default rate would a loss
  spread of 200 basis points cover?
  \item\label{ex6:eleven} A borrower has debt of 4 times EBITDA and a maintenance covenant at 5.5
  times. EBITDA has no drift and a volatility of 20\% a year. Use \cref{prop6:covenant} to find the
  probability of a breach within three years.
  \item\label{ex6:twelve} A fund buys back 5\% of its net assets each quarter, pro rata, and tenders
  total 12.5\% every quarter. What share of each request is filled, and what share of a holding is still
  in the fund after two years for an investor who tenders all of it every quarter?
  \item\label{ex6:thirteen} A BDC has debt of 1.25 times net assets and no other liabilities. What share of its assets can
  it lose before its asset coverage falls below 150\%, and what does a loss of 10\% of assets cost its
  shareholders?
  \item\label{ex6:fourteen} A borrower pays the base rate plus 6\% and has interest coverage of 3 when the base rate is
  1\%. What is its coverage, with the same EBITDA, when the base rate is 5\%?
\end{enumerate}

\subsection*{Computing projects}
\label{sec6:projects}
These use the data and code in the repository. Each is solved by a function in the file
\texttt{code/ch06/projects.py}, which the tests in \texttt{test\_projects.py} run, so an answer
can be checked by running that file from the root of the repository. The shipped series end in
September 2026, and those that end earlier are noted.

\begin{enumerate}[label=\textbf{P6.\arabic*}, leftmargin=*, ref=P6.\arabic*]
\item Rerun the refinancing case study with \texttt{refi\_vintages} for loan-to-value limits of 60\%, 65\% and
75\%, coverage tests of 1.20, 1.25 and 1.35, and cap rate sensitivities $\beta$ of 0, 0.5 and 1. Count the
vintages whose loan cannot be refinanced in full. Explain why, with $\beta$ fixed, the count barely moves
with the lending standards when the same standards apply at origination and at maturity, and why it moves
with $\beta$.
\item For each quoted maturity of \cref{tab6:hazard}, compute the upfront on a 100 basis point coupon from the
flat hazard that reprices the quote, with \texttt{flat\_hazard} and \texttt{cds\_upfront}, and from the
bootstrapped curve. Why do the two agree exactly at one year, and why does the gap widen with maturity on
this rising curve?
\item For each month since 2000, price a 10 year Government of Canada par bond at the month-end yield and a
month later, and compare the exact price change with the duration and the duration and convexity
approximations from \texttt{mod\_duration} and \texttt{convexity}. Report the root mean square error of each
in cents per 100 and the size and month of the largest miss of the duration approximation.
\item For each business development company in \texttt{data/bdc\_quarterly.csv}, read with
\texttt{load\_bdc} and \texttt{bdc\_panel} in \texttt{code/ch06/direct\_lending.py}, compute the first
autocorrelation of its quarterly return on net assets, its volatility before and after de-smoothing as
in Chapter~8, and the slope of its income yield on the one month bill rate. Which lenders look
most smoothed, and does the pass-through of the base rate differ much between them? The filings end in
June 2026.
\end{enumerate}

\subsection*{Answers}

\begin{enumerate}[label=\textbf{6.\arabic*}, leftmargin=*]
  \item The price is \xsOneP, the modified duration \xsOneD{} and the convexity \xsOneC. Duration alone
  gives \xsOneDur{} and duration with convexity \xsOneQd, against an exact price of \xsOneEx. The
  second estimate is slightly high, as \cref{prop6:shape} predicts for a rise in yield.
  \item The closed form and the sum both give \xsTwoD.
  \item The present value is $50e^{-0.24}=\xsThreePV$. Matching duration puts a share $(6-1)/(10-1)$ of
  the present value in the ten year bond, so the payments are \xsThreeA{} at one year and \xsThreeB{} at
  ten. By \cref{thm6:redington} the surplus is nonnegative after every parallel shift and zero only
  when there is no shift.
  \item By \cref{prop6:holee} it is $\tfrac12\times0.01^2\times3\times3.25$, which is \xsFour{} basis points.
  \item Equity is worth \xsFiveE{} and debt \xsFiveDebt, the spread is \xsFiveS{} basis points and the
  probability of default \xsFivePd\%.
  \item $\lambda=0.015/0.6=\xsSixL\%$ and $S(5)=e^{-5\lambda}=\xsSixS\%$.
  \item By \cref{prop6:lhp} the share is
  $\Phi\bigl((\Phi^{-1}(0.02)+\sqrt{0.2}\,\Phi^{-1}(0.99))/\sqrt{0.8}\bigr)=\xsSeven\%$.
  \item The loan to value is \xsEightLtv\%, the debt yield \xsEightDy\% and the coverage
  \xsEightDscr. The implied cap rate is 6\%, which equals $0.70\times\xsEightDy\%$, and the coverage is
  $\xsEightDy\%/7\%$.
  \item The balance due is \xsNineB. The new loan constant is \xsNineK\%, so the largest new loan is
  $90{,}000/(1.25\times\xsNineK\%)=\xsNineMax$ and the cover is \xsNineCover.
  \item By \cref{prop6:unitranche}(ii) the last out spread is $(600-0.4\times350)/0.6=\xdSl$ basis
  points. A loss spread of 2\% covers a default rate of $2\%/(1-0.33)=\xdPd\%$ a year.
  \item The barrier is $b=4/5.5$ of today's EBITDA. With no drift the probability is
  \[2\Phi\bigl(\ln(4/5.5)/(0.2\sqrt3)\bigr)=2\Phi(\xdA)=\xdBreach\%.\]
  \item $5/12.5=\xdFill\%$ of each request is filled, and after eight quarters $(1-0.4)^8=\xdLeft\%$ of
  the holding is still in the fund.
  \item By \cref{prop6:coverage} the fund can lose $1-1.5\times1.25/2.25=\xdThr\%$ of its assets, and a
  loss of 10\% of assets costs $2.25\times10\%=\xdEq\%$ of net assets.
  \item Coverage falls in proportion to $(1+6)/(5+6)$, from 3 to \xdIcr.
\end{enumerate}

% =====================================================================================
\section{Notes, sources and credits}
\label{sec6:credits}

\paragraph{Sources}
Immunisation is from \textcite{redington1952}, and its extension to coupon bonds over a holding period
from \textcite{fisher1971}. The short-rate models are from \textcite{vasicek1977}, \textcite{cox1985},
with the boundary behaviour from \textcite{feller1951}, and \textcite{ho1986}; \textcite{brigo2006}
treats all three, the forward measure and caplet pricing at length. The zero-coupon bond option
formula is from \textcite{jamshidian1989} and the caplet formula used for quoting from
\textcite{black1976}. The structural model is from \textcite{merton1974}. Reduced-form models are from
\textcite{jarrow1995}, \textcite{lando1998} and \textcite{duffie1999}, and \textcite{okane2008}
covers the conventions of credit default swaps and the pricing of tranches. The Gaussian copula for
default times is from \textcite{li2000} and the large-pool limit from \textcite{vasicek2002}.
\textcite{hull2022} covers market conventions for swaps, caps and credit derivatives. The description of
private credit and direct lending and the recoveries and interest coverage of direct loans are from
\textcite{cai2024}, the size of the market and its vulnerabilities from \textcite{imf2024gfsr}, and the
Canadian figures from \textcite{boc2026private} and \textcite{bocfsr2026}. The terms of a listed BDC are
from the annual report of Ares Capital \parencite{arcc2025} and those of a non-listed one from the annual
report of Blackstone Private Credit Fund \parencite{bcred2025}. The minimum of a Brownian motion with drift
is in \textcite{shreve2004}, the continuity correction for discrete monitoring is from
\textcite{broadie1997}, smoothing follows \textcite{geltner1993} as in Chapter~8, and the lagged beta is
from \textcite{dimson1979}. The unitranche split, the covenant recovery, the fee on gross assets and the
tender queue are short derivations given here.

\paragraph{Software used}
The scripts call NumPy \parencite{numpy} for arrays and random numbers, SciPy \parencite{scipy} for the normal
distribution, quadrature and root finding, pandas \parencite{pandas} to read the yield file, and
Matplotlib \parencite{matplotlib} for the figures. The cross-checks run in R \parencite{rcore} and
GNU Octave \parencite{octave}, and the screener's own calculation file runs in Node.js \parencite{nodejs}. The reference for
each package is the version used for the numbers printed here, and
\texttt{code/requirements-lock.txt} records the full environment.

\paragraph{Projects for comparison}
The private credit screener \texttt{private-credit-screener} \parencite{pcscreener} supplies the ratios, scoring rules and example
loans of \cref{sec6:private}; a copy of its calculation file, \texttt{calc.js}, is kept in
\texttt{code/ch06/vendor/} so that the tests can run it. The code here does not call the following
projects, but they implement the same methods in production form and are useful for checking:
QuantLib \parencite{quantlib} has short-rate models, caplet engines, CDS pricing and hazard curve bootstrapping;
FinancePy \parencite{financepy} has CDS, CDS index and tranche pricers in Python; the Open Source Risk Engine \parencite{ore}, built on
QuantLib, adds exposure and portfolio analytics; and gs-quant \parencite{gsquant} shows how a dealer's toolkit
organises rates and credit instruments. The full list, with links, is in \texttt{CREDITS.md}.

\paragraph{Data}
The case study uses the Government of Canada five year benchmark bond yield, series
\texttt{BD.CDN.5YR.DQ.YLD}, from the Bank of Canada Valet service \parencite{bocvalet}, in the file of
Chapter~4, which covers \bdFirst{} to \bdDate{} and was retrieved on 19 October 2026. The request
ended on 30 September 2026, so no October value enters any result, and the Bank's file has no value
for 30 September, so the last observation is \bdDate. The
study uses the last daily value of each month, so the last month-end is \csMatLast. The Bank can revise
published values, and the numbers here use the saved file. The accounts of the business development
companies in \cref{sec6:direct} are from the XBRL company facts of the SEC \parencite{secxbrl}, read on
22 October 2026 by \texttt{scripts/pull\_bdc.py} from the 10-Q and 10-K filings of \dlFirms{} companies
for quarters to June 2026; information on the SEC's site is public and may be redistributed, so the
table and the facts it is built from are shipped. The bill rate is the US one month series of the
Chapter~5 file. The credit fund returns are the Form PF statistics of the Office of Financial Research
used in Chapter~9 \parencite{ofrhfm}. The equity market return is from the French data library
\parencite{frenchlib}, which is not redistributed. The unitranche, the covenant example and the tender examples are
illustrations. The CDS quotes in \cref{sec6:hazard}, the
firm in \cref{sec6:merton} and the loans in \cref{sec6:private} are illustrations and describe no real
borrower.

\end{chapterrefs}
